\chapter{稳定性与李雅普诺夫方法（基础）}
\label{ch:lyapunov-basics}

% ════════════════════════════════════════════════════════════════
%  P4 控制部「C0 现代控制理论基础」子部 · 第 C0-3 章
%  主源：刘豹《现代控制理论》(第3版) 第4章（正文主线·物理直觉·几何图解·一题多解）
%  辅源/互校：张嗣瀛·高立群《现代控制理论》(现控02) 第5章
%             （克拉索夫斯基法的赫米特表述·用 V 求参数最优化·用 V 估计响应快速性·离散李氏）
%  密度梯度：导论 \cref{ch:control} 已"给会用"（def:ctrl-lyap-stable / thm:ctrl-lyap-direct /
%             thm:ctrl-lyap-eq / thm:ctrl-indirect 等），本章"给会推+物理直觉+会造 V"；
%             严格 ε-delta 证明、ISS、反步、逆定理 → \cref{ch:lyapunov-stability}（A–Q 的 D 章，做透层）。
%  记号：状态粗体 $\mathbf{x}$；李氏方程 $\mathbf{A}^\top\mathbf{P}+\mathbf{P}\mathbf{A}=-\mathbf{Q}$；
%        教材源非粗体/P/Q 已统一粗体化。控制侧 P/Q/R 与估计侧 Σ 双关纪律见首节 note。
%  OCR 勘误已逐处对照：刘豹导数点系统性丢失（\dot V↔V）、ex4-10 的 $\dot V=-x_3^2$ 负号、
%        Krasovskii $V=\dot{\mathbf x}^\top\mathbf P\dot{\mathbf x}$、$\mathbf x_{e3}=(0,1)$、
%        "设 $\mathbf A=\boldsymbol\Lambda$"对角化；现控02 旋度方程分母 $\partial x_1$（非 $\partial\dot x_1$）、
%        $\lim_{t\to\infty}$。详见结构化返回 ocrFixesApplied。
% ════════════════════════════════════════════════════════════════

\begin{practice}[前置自测]
开始之前，先试答下列问题。若已能流畅作答，可只速读本章的图解与速查卡；若卡住，正是本章要为你补的“地基”。
\begin{enumerate}
  \item “一个系统稳定”到底是什么意思？“受扰后留在附近”“最终回到平衡”“以多快回到”这三句话，分别对应哪一层稳定性？
  \item 给你一个非线性系统 $\dot{\mathbf{x}}=\mathbf{f}(\mathbf{x})$，你怎样\emph{不解微分方程}就判断它在原点是否渐近稳定？你打算用什么“量”来代替能量？
  \item 对线性系统 $\dot{\mathbf{x}}=\mathbf{A}\mathbf{x}$，要判稳定，你会去解特征值，还是去解一个矩阵方程 $\mathbf{A}^\top\mathbf{P}+\mathbf{P}\mathbf{A}=-\mathbf{Q}$？这两条路为什么等价？
  \item 李雅普诺夫第二法的“命门”是构造一个李雅普诺夫函数 $V(\mathbf{x})$。除了“凑二次型 $\mathbf{x}^\top\mathbf{P}\mathbf{x}$”，你还知道哪些\emph{可操作}的造 $V$ 流程？
  \item 同一个系统，选不同的 $V$ 会得到不同的“渐近稳定区域”吗？哪个 $V$ 算“更好”？
\end{enumerate}
\end{practice}

\section*{本章目标}
学完本章，你应当能够：
\begin{enumerate}
  \item \textbf{说清}李雅普诺夫意义下稳定、渐近稳定、大范围渐近稳定、不稳定四种定义，并用 $\varepsilon$-$\delta$ 球域图把它们\emph{画}出来；
  \item \textbf{区分}李雅普诺夫第一法（间接法/线性化）与第二法（直接法/能量法），知道前者何时\emph{失效}、为何需要后者；
  \item \textbf{建立} $V$ 是“到平衡的广义距离/能量”、$\dot V$ 是“趋近速度”的物理图像，并据此理解四条稳定性判据；
  \item \textbf{推导并会算}线性系统的李雅普诺夫方程 $\mathbf{A}^\top\mathbf{P}+\mathbf{P}\mathbf{A}=-\mathbf{Q}$——给定 $\mathbf{Q}$ 解 $\mathbf{P}$、用希尔维斯特判据判正定，从而判稳定、定参数范围；
  \item \textbf{掌握}两种“如何造 $V$”的可操作流程——克拉索夫斯基（雅可比）法与变量梯度法，并能在同一系统上\emph{一题多解}比较所得 $V$ 的优劣；
  \item \textbf{运用} $V$ 的两项进阶用途——求参数最优化、估计动态响应的快速性（衰减率/时间常数）；
  \item \textbf{衔接}：把本章建立的“会用、会推、会造 $V$”交给 \cref{ch:lyapunov-stability} 去做严格 $\varepsilon$-$\delta$ 证明、ISS、反步与逆定理，并理解 \cref{ch:dp-hjb} 的 HJB 值函数为何也是一个李雅普诺夫函数。
\end{enumerate}

\subsection*{本章在“现控基础”主线中的位置}
本章是 P4 控制部「现代控制理论基础」子部（C0）的第三站，上承状态空间建模（\cref{ch:statespace-basics}）与能控能观（\cref{ch:controllability-obsv}），下接综合与最优控制基础（\cref{ch:linear-synthesis,ch:optimal-control-basics}）。一条主线是：\emph{先把“稳定”这个最核心的系统属性讲透——从直觉到定义、从判据到“怎么造判稳定的工具 $V$”}，因为后面一切收敛性论证（极点配置闭环、观测器误差、LQR 闭环、PD+重力补偿、MPC 的终端代价）都以它为统一语言。

\subsection*{与控制导论的分工（务必先读）}
导论 \cref{ch:control} 已在 \cref{sec:ctrl-lyapunov} 给出了\emph{会用}层面的全部结论：三层定义（\cref{def:ctrl-lyap-stable,def:ctrl-asymp-stable,def:ctrl-exp-stable}）、直接法（\cref{thm:ctrl-lyap-direct}）、LaSalle 不变原理（\cref{thm:ctrl-lasalle}）、线性系统的李雅普诺夫方程（\cref{thm:ctrl-lyap-eq}）、间接法（\cref{thm:ctrl-indirect}）与 Barbalat 引理（\cref{lem:ctrl-barbalat}）。\textbf{本章不重复编号这些定义与定理}，而是\emph{把它们展开}：补上导论“点到为止”的物理直觉（能量小球、$\varepsilon$-$\delta$ 球域、相平面水平集）、把线性李氏方程\emph{亲手推一遍并算几个例子}、并系统补全导论\emph{完全没讲}的两件大事——“\textbf{怎样构造 $V$}”（克拉索夫斯基法、变量梯度法）与“\textbf{$V$ 还能干什么}”（求参数最优化、估计响应速度）。读到导论已陈述处，本章用半句交叉引用接住即下沉。

\begin{note}
\textbf{本章记号与约定.}\;
状态向量粗体 $\mathbf{x}\in\mathbb{R}^n$，自治系统 $\dot{\mathbf{x}}=\mathbf{f}(\mathbf{x})$（或时变 $\dot{\mathbf{x}}=\mathbf{f}(\mathbf{x},t)$），平衡状态 $\mathbf{x}_e$。
李雅普诺夫函数记 $V(\mathbf{x})$（标量），其沿轨迹的时间导数记 $\dot V(\mathbf{x})$。
二次型李氏函数 $V(\mathbf{x})=\mathbf{x}^\top\mathbf{P}\mathbf{x}$，$\mathbf{P}=\mathbf{P}^\top$；线性李氏方程统一写作 $\mathbf{A}^\top\mathbf{P}+\mathbf{P}\mathbf{A}=-\mathbf{Q}$（与导论 \cref{eq:ctrl-lyap-eq} 同形，引用即可）。
\textbf{两本教材源}用非粗体、且常把矩阵记作 $P,Q$；融入时一律粗体化为 $\mathbf{P},\mathbf{Q}$。
\textbf{P/Q/R 双关纪律}：本章及控制侧 $\mathbf{P}$=李氏函数权重/代价权重、$\mathbf{Q},\mathbf{R}$=代价权重；估计侧（\cref{ch:eskf}）噪声协方差另记 $\boldsymbol\Sigma$、Kalman 增益 $\mathbf{L}$，两者\emph{不要混用同一字母}。
矩阵不等式 $\mathbf{P}\succ0$（正定）、$\mathbf{P}\succeq0$（半正定），与教材“记作 $\mathbf{P}>0$”同义。
\end{note}

\subsection*{本章知识导航}
本章是一条“\emph{从‘什么叫稳定’到‘怎样造出判稳定的工具’}”的主线。结构如下：

\begin{center}
\small
\begin{tikzpicture}[
  node distance=4.5mm and 7mm, >=Stealth,
  blk/.style={draw, rounded corners, align=center, font=\small, inner sep=3pt, fill=main!6, draw=main!60!black},
  grp/.style={draw, rounded corners, align=center, font=\small, inner sep=3pt, fill=second!6, draw=second!60!black},
  toolb/.style={draw, rounded corners, align=center, font=\small, inner sep=3pt, fill=third!8, draw=third!60!black},
  ln/.style={->, thick, gray!60!black}]
\node[blk] (def) {稳定性定义\\ $\varepsilon$-$\delta$ 球域图};
\node[blk, right=of def] (first) {第一法\\（线性化）};
\node[grp, right=of first] (second) {第二法\\（能量·四判据）};
\node[grp, below=8mm of def] (prep) {造 $V$ 预备\\ 定号·二次型·Sylvester};
\node[grp, right=of prep] (lineq) {线性李氏方程\\ $\mathbf{A}^\top\mathbf{P}+\mathbf{P}\mathbf{A}=-\mathbf{Q}$};
\node[blk, right=of lineq] (disc) {时变 / 离散\\ 李氏判据};
\node[toolb, below=8mm of prep] (kras) {克拉索夫斯基法\\（造 $V$ 流程一）};
\node[toolb, right=of kras] (vargrad) {变量梯度法\\（造 $V$ 流程二）};
\node[toolb, right=of vargrad] (use) {$V$ 进阶用途\\ 参数最优·响应速度};
\draw[ln] (def)--(first); \draw[ln] (first)--(second);
\draw[ln] (second.south) -- ++(0,-4.5mm) -| (lineq.north);
\draw[ln] (prep)--(lineq); \draw[ln] (lineq)--(disc);
\draw[ln] (prep.south)--(kras.north);
\draw[ln] (kras)--(vargrad); \draw[ln] (vargrad)--(use);
\end{tikzpicture}
\end{center}

\noindent\textbf{推荐路径}：第 1--3 节是任何读者的主干（定义 + 两法 + 四判据 + 造 $V$ 预备）；第 4 节把线性李氏方程\emph{推透会算}（与 LQR/极点配置闭环分析直接相关，\cref{ch:linear-synthesis,ch:optimal-control-basics} 反复引用）；第 5 节的两种“造 $V$ 流程”是\emph{差异化干货}（导论完全没讲）；第 6 节的 $V$ 进阶用途面向控制综合，可在读完最优控制基础后回看。带（互校）标记处给出刘豹与现控02 的独立第二信源对照。

% ════════════════════════════════════════════════════════════════
\section{李雅普诺夫意义下的稳定性定义}
\label{sec:lyap-def}

\paragraph{动机：为什么经典稳定判据不够用。}
经典控制理论里判稳定有现成、好用的工具：单输入单输出线性定常系统用劳斯（Routh）、胡维茨（Hurwitz）代数判据，频域用奈奎斯特（Nyquist）判据\cite{khalil2002nonlinear}。它们都以“分析特征方程的根在复平面上的分布”为基础。可一旦面对\emph{非线性系统}或\emph{时变系统}，特征方程根的概念就失效了——非线性系统甚至连“线性系统的稳定性只取决于结构与参数、与初始条件无关”这条经验都不再成立：非线性系统的稳定性\emph{还与初始条件和扰动的大小有关}。这正是 1892 年俄国数学家李雅普诺夫（A.\,M.~Lyapunov）在《运动稳定性的一般问题》中要解决的：给出一个\emph{对线性、非线性、定常、时变系统都普遍适用}的稳定性一般定义\cite{lyapunov1992general}。他把判稳定的方法归为两类——第一法（间接法，解方程看解的性质）与第二法（直接法，不解方程、用一个标量函数判定）。本节先立定义，后两节分别展开两法。

\paragraph{平衡状态：一切稳定性都是相对它而言的。}
设自治系统的齐次状态方程为
\begin{equation}
  \dot{\mathbf{x}}=\mathbf{f}(\mathbf{x},t),
  \label{eq:lyap-sys}
\end{equation}
其中 $\mathbf{x}\in\mathbb{R}^n$、$\mathbf{f}$ 是与 $\mathbf{x}$ 同维的向量函数。若 $\mathbf{f}$ 不显含 $t$ 即为定常（自治）系统。给定初值 $(t_0,\mathbf{x}_0)$ 的唯一解记为 $\mathbf{x}=\boldsymbol\Phi(t;\mathbf{x}_0,t_0)$，它在 $n$ 维状态空间中画出一条从 $\mathbf{x}_0$ 出发的\emph{状态轨线}。若存在状态 $\mathbf{x}_e$ 对所有 $t$ 都使
\begin{equation}
  \mathbf{f}(\mathbf{x}_e,t)\equiv\mathbf{0},
  \label{eq:equilibrium}
\end{equation}
则称 $\mathbf{x}_e$ 为\textbf{平衡状态}（平衡点）。

\begin{insight}
平衡点不一定唯一，也不一定存在。线性定常系统 $\dot{\mathbf{x}}=\mathbf{A}\mathbf{x}$：$\mathbf{A}$ 非奇异时 $\mathbf{x}_e=\mathbf{0}$ 是\emph{唯一}平衡点；$\mathbf{A}$ 奇异时有\emph{无穷多个}平衡点。非线性系统通常有一个或多个孤立平衡点，由代数方程 \cref{eq:equilibrium} 的常值解给出。\textbf{关键约定}：任何一个已知平衡点都可经坐标平移 $\mathbf{y}=\mathbf{x}-\mathbf{x}_e$ 搬到原点，故\emph{今后只讨论原点 $\mathbf{x}_e=\mathbf{0}$ 的稳定性}。这也解释了为什么“线性定常系统笼统讲系统稳定性”是合法的（只有唯一平衡点），而非线性系统必须\emph{逐个平衡点}分别讨论——不同平衡点可能稳定性截然不同。
\end{insight}

\noindent\textbf{多平衡点示例.}\;系统 $\dot x_1=-x_1$、$\dot x_2=x_1+x_2-x_2^3$ 有三个平衡点 $\mathbf{x}_{e1}=(0,0)^\top$、$\mathbf{x}_{e2}=(0,-1)^\top$、$\mathbf{x}_{e3}=(0,1)^\top$。% OCR: 刘豹源第三个误写 x_{e1}，据"三个平衡状态"应为 x_{e3}
它们各自的稳定性需分别判断——这正是非线性稳定性分析与经典理论的根本不同。

\subsection{四种稳定性定义与球域图}
\label{sec:lyap-def-four}

李雅普诺夫稳定性的精髓在于用两个超球域刻画“受扰后跑不跑得远、回不回得来”。用 $\|\mathbf{x}-\mathbf{x}_e\|$（欧几里得范数）表示状态到平衡点的距离，以 $\mathbf{x}_e$ 为心、半径 $\varepsilon$ 的超球域记 $S(\varepsilon)$，则 $\mathbf{x}\in S(\varepsilon)$ 意味 $\|\mathbf{x}-\mathbf{x}_e\|\le\varepsilon$。下面四种定义可与导论 \cref{def:ctrl-lyap-stable,def:ctrl-asymp-stable} 的紧凑形式对照阅读——导论给定义，本章给\emph{图与边界}。

\begin{definition}[李雅普诺夫意义下稳定（球域形式）]\label{def:lyap-isl}
平衡状态 $\mathbf{x}_e$ 称\textbf{李雅普诺夫意义下稳定}，若对任意给定实数 $\varepsilon>0$，都存在另一实数 $\delta(\varepsilon,t_0)>0$，使当 $\|\mathbf{x}_0-\mathbf{x}_e\|\le\delta$ 时，从任意初态 $\mathbf{x}_0$ 出发的解满足
\begin{equation}
  \|\boldsymbol\Phi(t;\mathbf{x}_0,t_0)-\mathbf{x}_e\|\le\varepsilon,\qquad t_0\le t<\infty.
  \label{eq:lyap-isl}
\end{equation}
若 $\delta$ 可取得与 $t_0$ 无关，则进一步称\textbf{一致稳定}。
\end{definition}

\begin{insight}
读法是“\emph{先有 $\varepsilon$，后找 $\delta$}”：外圈 $S(\varepsilon)$ 是你\emph{要求}解永不越出的边界（任意小）；内圈 $S(\delta)$ 是\emph{找得到}的一个初值容差。只要“对每个 $\varepsilon$ 都配得出一个 $\delta$，使从 $S(\delta)$ 出发的轨线永远困在 $S(\varepsilon)$ 内”，就叫稳定。$S(\delta)$ 限制初态、$S(\varepsilon)$ 规定自由响应的边界——一言以蔽之：\textbf{初态够近 $\Rightarrow$ 全程不跑远}。注意它\emph{不}要求最终回到 $\mathbf{x}_e$。
\end{insight}

\begin{definition}[渐近稳定 / 大范围渐近稳定 / 不稳定]\label{def:lyap-asymp}
设 $\mathbf{x}_e$ 已是李雅普诺夫意义下稳定的。
\begin{enumerate}
  \item \textbf{渐近稳定}：若当 $t\to\infty$ 时轨线不仅不超出 $S(\varepsilon)$，而且\emph{最终收敛于} $\mathbf{x}_e$，即 $\lim_{t\to\infty}\|\boldsymbol\Phi(t;\mathbf{x}_0,t_0)-\mathbf{x}_e\|=0$；
  \item \textbf{大范围（全局）渐近稳定}：若从状态空间\emph{所有}初态出发的轨线都渐近稳定（即 $\delta$ 可取 $\infty$）。其\emph{必要条件}是整个状态空间只有一个平衡点；
  \item \textbf{不稳定}：若存在某个 $\varepsilon>0$，使得无论 $\delta>0$ 取多小，从 $S(\delta)$ 内出发的轨线中\emph{至少有一条}越出 $S(\varepsilon)$。
\end{enumerate}
\end{definition}

\begin{figure}[H]
\centering
\begin{tikzpicture}[>=Stealth, font=\small]
% ---- (a) Lyapunov stable: stays in S(eps) but not converging ----
\begin{scope}
  \fill[main!7] (0,0) circle (1.55);
  \fill[second!12] (0,0) circle (0.62);
  \draw[main!70!black, thick] (0,0) circle (1.55);
  \draw[second!70!black, thick, dashed] (0,0) circle (0.62);
  \fill (0,0) circle (1.1pt) node[below left=0pt and 1pt, font=\scriptsize] {$\mathbf{x}_e$};
  \node[font=\scriptsize, main!60!black] at (1.15,1.25) {$S(\varepsilon)$};
  \node[font=\scriptsize, second!60!black] at (0.0,0.86) {$S(\delta)$};
  \fill[third!70!black] (0.45,0.30) circle (0.9pt) node[right=0pt, font=\scriptsize]{$\mathbf{x}_0$};
  % a meandering bounded trajectory that never settles
  \draw[third, thick, ->] (0.45,0.30)
        .. controls (1.3,0.6) and (1.0,-1.1) .. (-0.2,-1.2)
        .. controls (-1.35,-1.3) and (-1.4,0.6) .. (-0.3,1.15)
        .. controls (0.7,1.5) and (1.45,0.7) .. (1.2,-0.2);
  \node[align=center, font=\scriptsize] at (0,-2.25) {(a) 李雅普诺夫意义下稳定\\轨线困在 $S(\varepsilon)$ 内，\emph{未必}回到 $\mathbf{x}_e$};
\end{scope}
% ---- (b) asymptotically stable: spirals into x_e ----
\begin{scope}[xshift=5.0cm]
  \fill[main!7] (0,0) circle (1.55);
  \fill[second!12] (0,0) circle (0.62);
  \draw[main!70!black, thick] (0,0) circle (1.55);
  \draw[second!70!black, thick, dashed] (0,0) circle (0.62);
  \fill (0,0) circle (1.1pt) node[above right=0pt and 1pt, font=\scriptsize]{$\mathbf{x}_e$};
  \fill[third!70!black] (0.5,0.28) circle (0.9pt) node[right=0pt, font=\scriptsize]{$\mathbf{x}_0$};
  \draw[third, thick, ->] (0.5,0.28)
        .. controls (1.0,0.8) and (0.2,1.15) .. (-0.55,0.55)
        .. controls (-1.0,0.1) and (-0.6,-0.6) .. (-0.05,-0.55)
        .. controls (0.45,-0.5) and (0.5,-0.05) .. (0.18,0.2)
        .. controls (0.05,0.32) and (-0.05,0.1) .. (0,0);
  \node[align=center, font=\scriptsize] at (0,-2.25) {(b) 渐近稳定\\轨线收敛到 $\mathbf{x}_e$};
\end{scope}
% ---- (c) unstable: leaves S(eps) ----
\begin{scope}[xshift=10.0cm]
  \fill[main!7] (0,0) circle (1.55);
  \fill[second!12] (0,0) circle (0.62);
  \draw[main!70!black, thick] (0,0) circle (1.55);
  \draw[second!70!black, thick, dashed] (0,0) circle (0.62);
  \fill (0,0) circle (1.1pt) node[below left=0pt and 1pt, font=\scriptsize]{$\mathbf{x}_e$};
  \fill[third!70!black] (0.18,0.12) circle (0.9pt) node[below right=-1pt and 0pt, font=\scriptsize]{$\mathbf{x}_0$};
  \draw[third, thick, ->] (0.18,0.12)
        .. controls (0.45,0.6) and (-0.35,0.85) .. (-0.55,0.25)
        .. controls (-0.85,-0.5) and (0.45,-0.9) .. (0.9,-0.35)
        .. controls (1.5,0.35) and (0.6,1.4) .. (-0.55,1.75);
  \node[align=center, font=\scriptsize] at (0,-2.25) {(c) 不稳定\\至少一条轨线越出 $S(\varepsilon)$};
\end{scope}
\end{tikzpicture}
\caption{李雅普诺夫稳定性的 $\varepsilon$-$\delta$ 球域几何（二阶系统，仿刘豹\cite{liubao2006modern} 图 4.1--4.3 重绘）。外圈实线为 $S(\varepsilon)$（要求轨线不越出），内虚线圈为 $S(\delta)$（初值容差）。(a) 稳定但不收敛——经典理论称之为临界稳定；(b) 渐近稳定——既不跑远又最终回到平衡；(c) 不稳定。}
\label{fig:lyap-ball-domain}
\end{figure}

\begin{note}
\textbf{两种“稳定”的口径差异（与经典理论的接口）.}\;
简言之：若自由响应 $\mathbf{x}(t)=\boldsymbol\Phi(t;\mathbf{x}_0,t_0)$ \emph{有界}则 $\mathbf{x}_e$ 稳定；若有界\emph{且} $\lim_{t\to\infty}\mathbf{x}(t)=\mathbf{0}$ 则\emph{渐近}稳定；若无界则\emph{不稳定}。\textbf{要警惕}：在经典控制理论里“稳定系统”专指\emph{渐近}稳定；仅在李雅普诺夫意义下稳定但非渐近稳定的（如无阻尼振子，\cref{fig:lyap-ball-domain}a），经典理论称\emph{临界稳定}，工程上归入不稳定。本书行文沿用李雅普诺夫口径，凡说“稳定”均指 \cref{def:lyap-isl}，需要更强结论时显式写“渐近”。现控02\cite{zhang2017modern} 的定义与此完全一致，可作独立第二信源互校（互校）。
\end{note}

% ════════════════════════════════════════════════════════════════
\section{李雅普诺夫第一法（间接法）}
\label{sec:lyap-first}

第一法的基本思路最朴素：\emph{解出系统的运动方程，再看解的性质}——这与经典理论一脉相承。对线性定常系统，只需解特征方程的根；对非线性系统，则先\emph{线性化}取一次近似，再看线性化矩阵的特征根。

\subsection{线性系统的稳定判据}
\label{sec:lyap-first-linear}

\begin{theorem}[线性定常系统稳定判据]\label{thm:lyap-linear-eig}
线性定常系统 $\dot{\mathbf{x}}=\mathbf{A}\mathbf{x}+\mathbf{b}u,\ y=\mathbf{c}\mathbf{x}$ 的平衡状态 $\mathbf{x}_e=\mathbf{0}$ \textbf{渐近稳定的充要条件}是矩阵 $\mathbf{A}$ 的所有特征值均具有负实部（即 $\mathbf{A}$ 为 Hurwitz 阵）。
\end{theorem}

\noindent 这与导论 \cref{thm:ctrl-lyap-eq} 的第一句一致。导论从 LQR 闭环分析的角度陈述它；这里强调它是\emph{状态（内部）稳定性}的判据。工程上还常关心\emph{输出稳定性}：

\begin{proposition}[输出稳定与零极点对消]\label{prop:output-stable}
线性定常系统 $\Sigma:(\mathbf{A},\mathbf{b},\mathbf{c})$ \textbf{输出稳定}（有界输入引起有界输出）的充要条件是其传递函数 $W(s)=\mathbf{c}(s\mathbf{I}-\mathbf{A})^{-1}\mathbf{b}$ 的极点全部位于 $s$ 左半平面。仅当 $W(s)$ \emph{无零极点对消}（$\mathbf{A}$ 的特征值与 $W(s)$ 的极点相同）时，状态稳定性才与输出稳定性一致。
\end{proposition}

\begin{insight}[一个会“骗过”输出的不稳定系统]
考虑
\[
  \dot{\mathbf{x}}=\begin{bmatrix}-1&0\\0&1\end{bmatrix}\mathbf{x}+\begin{bmatrix}1\\1\end{bmatrix}u,\qquad y=(1,\ 0)\mathbf{x}.
\]
特征值 $\lambda_1=-1,\ \lambda_2=+1$，故\emph{状态不渐近稳定}。但传递函数
\[
  W(s)=(1,0)\begin{bmatrix}s+1&0\\0&s-1\end{bmatrix}^{-1}\!\begin{bmatrix}1\\1\end{bmatrix}=\frac{s-1}{(s+1)(s-1)}=\frac{1}{s+1},
\]
极点 $s=-1$ 在左半平面，故\emph{输出稳定}。原因是正实部特征值 $\lambda_2=+1$ 被零点 $s=+1$ \emph{对消}了，在输入输出特性里“看不见”。这正是为什么只看传递函数会漏判内部不稳定——\textbf{能控能观与稳定性必须合看}（对消即不能控或不能观，见 \cref{ch:controllability-obsv}）。
\end{insight}

\subsection{非线性系统的线性化判据}
\label{sec:lyap-first-nonlinear}

设非线性系统 $\dot{\mathbf{x}}=\mathbf{f}(\mathbf{x},t)$，$\mathbf{x}_e$ 为平衡点，$\mathbf{f}$ 对 $\mathbf{x}$ 有连续偏导。在 $\mathbf{x}_e$ 邻域把 $\mathbf{f}$ 展成泰勒级数：
\begin{equation}
  \dot{\mathbf{x}}=\frac{\partial\mathbf{f}}{\partial\mathbf{x}}\Big|_{\mathbf{x}_e}(\mathbf{x}-\mathbf{x}_e)+\mathbf{R}(\mathbf{x}),
  \label{eq:taylor-lin}
\end{equation}
其中雅可比矩阵 $\dfrac{\partial\mathbf{f}}{\partial\mathbf{x}}=\big[\partial f_i/\partial x_j\big]_{n\times n}$，$\mathbf{R}(\mathbf{x})$ 是高阶项。令 $\Delta\mathbf{x}=\mathbf{x}-\mathbf{x}_e$、取一次近似得线性化方程
\begin{equation}
  \Delta\dot{\mathbf{x}}=\mathbf{A}\,\Delta\mathbf{x},\qquad \mathbf{A}=\frac{\partial\mathbf{f}}{\partial\mathbf{x}}\Big|_{\mathbf{x}=\mathbf{x}_e}.
  \label{eq:lin-A}
\end{equation}

\begin{theorem}[李雅普诺夫间接法]\label{thm:lyap-indirect}
对 \cref{eq:lin-A} 的系数矩阵 $\mathbf{A}$（即 $\mathbf{f}$ 在 $\mathbf{x}_e$ 处的雅可比）：
\begin{enumerate}
  \item 若 $\mathbf{A}$ 的所有特征值实部 $<0$，则原非线性系统在 $\mathbf{x}_e$ \textbf{渐近稳定}，且稳定性与 $\mathbf{R}(\mathbf{x})$ 无关；
  \item 若 $\mathbf{A}$ 有特征值实部 $>0$，则 $\mathbf{x}_e$ \textbf{不稳定}；
  \item 若 $\mathbf{A}$ 有特征值实部 $=0$（临界情形），则稳定性\emph{取决于高阶项 $\mathbf{R}(\mathbf{x})$}，\textbf{不能}由 $\mathbf{A}$ 的特征值判定——间接法\emph{失效}。
\end{enumerate}
\end{theorem}

\noindent 这正是导论 \cref{thm:ctrl-indirect} 的展开版。第 3 条的“临界失效”是第一法最大的局限，也是第二法登场的理由。

\begin{example}[间接法的成功与失效各一例]\label{ex:lyap-indirect}
系统 $\dot x_1=x_1-x_1x_2$、$\dot x_2=-x_2+x_1x_2$ 有两个平衡点 $\mathbf{x}_{e1}=(0,0)^\top$、$\mathbf{x}_{e2}=(1,1)^\top$。
\begin{itemize}
  \item 在 $\mathbf{x}_{e1}$ 处线性化得 $\mathbf{A}=\big[\begin{smallmatrix}1&0\\0&-1\end{smallmatrix}\big]$，特征值 $\lambda=\pm1$，含正实部 $\Rightarrow$ 原系统在 $\mathbf{x}_{e1}$ \emph{不稳定}（第一法即可判定）。
  \item 在 $\mathbf{x}_{e2}$ 处线性化得 $\mathbf{A}=\big[\begin{smallmatrix}0&-1\\1&0\end{smallmatrix}\big]$，特征值 $\pm\mathrm{j}1$，实部为零 $\Rightarrow$ 临界情形，\emph{第一法无能为力}，必须用下一节的第二法。
\end{itemize}
\end{example}

\begin{pitfall}[思维陷阱：以为“线性化稳定 = 原系统稳定”总成立]
间接法只在特征值\emph{严格}落在左/右半平面时给出确定结论，且只是\emph{局部}的。临界情形（纯虚特征值）下，任意小的高阶项都可能把系统推向稳定或不稳定——典型如无阻尼振子线性化后是纯虚根。把“线性化矩阵临界”草率读成“原系统临界稳定”是常见错误。此时唯一出路是第二法（直接法）或 LaSalle（\cref{thm:ctrl-lasalle}）。
\end{pitfall}

% ════════════════════════════════════════════════════════════════
\section{李雅普诺夫第二法（直接法）}
\label{sec:lyap-second}

\paragraph{动机与能量直觉。}
第二法（直接法）的天才在于：\emph{不求解运动方程}，而借助一个标量函数直接判定稳定性。它从\emph{能量}的观点出发。一个受激励的物理系统，若其储存的能量随时间\emph{单调衰减}、到平衡时取最小，则平衡渐近稳定；若不断从外界吸能、储能越来越大，则不稳定；若储能既不增也不减，则为李雅普诺夫意义下的稳定。

\begin{figure}[H]
\centering
\begin{tikzpicture}[>=Stealth, font=\small, scale=1.0]
% energy bowl
\draw[thick, main!70!black] (-2.6,2.0) .. controls (-1.2,-0.55) and (1.2,-0.55) .. (2.6,2.0);
\node[main!60!black, font=\scriptsize] at (2.45,1.75) {$V(\mathbf{x})$};
% equilibrium A at bottom
\fill (0,-0.30) circle (1.4pt);
\node[font=\scriptsize, below=0pt] at (0,-0.34) {$A$ (平衡 $\mathbf{x}_e$)};
% perturbed point C up on the bowl
\fill[third!70!black] (1.75,1.05) circle (1.6pt);
\node[third!60!black, font=\scriptsize, right=0pt] at (1.75,1.05) {$C$ (受扰)};
% ball at C
\draw[third, thick, ->] (1.55,0.95) .. controls (0.7,0.2) and (-0.6,0.1) .. (-1.15,0.55);
\draw[third, thick, ->] (-1.0,0.45) .. controls (-0.4,0.0) and (0.5,0.05) .. (0.9,0.35);
% arrows for dissipation note
\node[align=left, font=\scriptsize] at (4.7,1.4)
  {有摩擦：$\dot V<0$\\储能衰减\\$\Rightarrow$ 渐近稳定};
\node[align=left, font=\scriptsize] at (4.7,-0.25)
  {无摩擦：$\dot V=0$\\等幅振荡\\$\Rightarrow$ 仅稳定};
\draw[->, gray!60!black] (3.0,1.4)--(2.3,1.1);
\draw[->, gray!60!black] (3.0,-0.25)--(2.4,0.1);
\end{tikzpicture}
\caption{能量法的物理图像（仿刘豹\cite{liubao2006modern} 图 4.4 重绘）。曲面上的小球受扰后从 $A$ 升到 $C$，获得能量后绕 $A$ 振荡。若曲面有摩擦，振荡过程耗能（$\dot V<0$），幅值越来越小、回到 $A$——渐近稳定；若曲面绝对光滑（$\dot V=0$），等幅振荡永不停——李雅普诺夫意义下稳定。$V$ 就是这个“广义能量”。}
\label{fig:energy-bowl}
\end{figure}

\noindent 难点在于：一般系统未必能直观写出能量函数。李雅普诺夫的办法是\emph{虚构}一个正定标量函数 $V(\mathbf{x})$ 当作广义能量，再看 $\dot V(\mathbf{x})=\mathrm{d}V/\mathrm{d}t$ 的符号——若 $V$ 正定而 $\dot V$ 负定，系统渐近稳定。这个 $V$ 就叫\textbf{李雅普诺夫函数}。任何满足判据条件的标量函数都可充任。于是\emph{用第二法的全部关键，归结为“找到一个合适的 $V$”}——这是本章后半部分（第 5 节）要专门攻克的难题。

\subsection{造 \texorpdfstring{$V$}{V} 的预备：标量函数的定号性}
\label{sec:lyap-definite}

要判 $V$ 与 $\dot V$ 的符号，先要把“正定/负定”说清楚。设 $V(\mathbf{x})$ 在域 $\Omega$ 上定义、$V(\mathbf{0})=0$，对 $\Omega$ 中任意非零 $\mathbf{x}$：

\begin{definition}[标量函数的定号性]\label{def:definiteness}
\begin{enumerate}
  \item $V(\mathbf{x})>0$ $\Rightarrow$ \textbf{正定}（如 $V=x_1^2+x_2^2$）；
  \item $V(\mathbf{x})\ge0$ $\Rightarrow$ \textbf{半正定（非负定）}（如 $V=(x_1+x_2)^2$，因 $\mathbf{x}=(a,-a)^\top$ 时为零）；
  \item $V(\mathbf{x})<0$ $\Rightarrow$ \textbf{负定}（如 $V=-(x_1^2+2x_2^2)$）；
  \item $V(\mathbf{x})\le0$ $\Rightarrow$ \textbf{半负定（非正定）}；
  \item $V$ 在 $\Omega$ 内既可正又可负 $\Rightarrow$ \textbf{不定}（如 $V=x_1+x_2$）。
\end{enumerate}
\end{definition}

\paragraph{二次型与希尔维斯特判据。}
李氏函数最常用的形式是\emph{二次型} $V(\mathbf{x})=\mathbf{x}^\top\mathbf{P}\mathbf{x}$，$\mathbf{P}=\mathbf{P}^\top$（实对称）。对实对称 $\mathbf{P}$ 必存在正交变换 $\mathbf{x}=\mathbf{T}\bar{\mathbf{x}}$ 使
\begin{equation}
  V(\mathbf{x})=\bar{\mathbf{x}}^\top\big(\mathbf{T}^{-1}\mathbf{P}\mathbf{T}\big)\bar{\mathbf{x}}=\sum_{i=1}^n\lambda_i\,\bar x_i^2,
  \label{eq:quad-canonical}
\end{equation}
其中 $\lambda_i$ 是 $\mathbf{P}$ 的（实）特征值。故 $V$ 正定 $\iff$ $\mathbf{P}$ 所有特征值 $>0$。$\mathbf{P}$ 的符号性质与 $V$ 完全一致，记 $\mathbf{P}\succ0$（正定）等。判 $\mathbf{P}$ 定号最实用的是\textbf{希尔维斯特（Sylvester）判据}：

\begin{theorem}[希尔维斯特判据]\label{thm:sylvester}
设实对称阵 $\mathbf{P}=[p_{ij}]_{n\times n}$，其各阶顺序主子式 $\Delta_1=p_{11}$、$\Delta_2=\big|\begin{smallmatrix}p_{11}&p_{12}\\p_{21}&p_{22}\end{smallmatrix}\big|$、$\dots$、$\Delta_n=\det\mathbf{P}$。则
\begin{enumerate}
  \item $\mathbf{P}$ 正定 $\iff$ 所有 $\Delta_i>0$；
  \item $\mathbf{P}$ 负定 $\iff$ $\Delta_i>0$（$i$ 偶）、$\Delta_i<0$（$i$ 奇），即\emph{符号交替、奇负偶正}；
  \item $\mathbf{P}$ 半正定 $\iff$ $\Delta_i\ge0\ (i<n)$ 且 $\Delta_n=0$。
\end{enumerate}
\end{theorem}

\begin{pitfall}[思维陷阱：把“对角元同号”当成定号判据]
$V=x_1^2+2x_1x_2+x_2^2+x_3^2$ 的矩阵 $\mathbf{P}=\big[\begin{smallmatrix}1&1&0\\1&1&0\\0&0&1\end{smallmatrix}\big]$ 对角元全正，却\emph{不是}正定——它是\emph{半}正定的（$\Delta_2=0$）。判定号性必须算\emph{顺序主子式}（\cref{thm:sylvester}）或特征值，不能只看对角元。这是手算李氏函数时最易翻车处。
\end{pitfall}

\subsection{四条稳定性判据}
\label{sec:lyap-criteria}

把能量直觉形式化。设 $\dot{\mathbf{x}}=\mathbf{f}(\mathbf{x})$，平衡点 $\mathbf{x}_e=\mathbf{0}$。下述判据是导论 \cref{thm:ctrl-lyap-direct} 的逐条展开（导论合并陈述，这里拆成“稳定/渐近/大范围/不稳定”四条便于对照例题）。

\begin{theorem}[李雅普诺夫第二法的四条判据]\label{thm:lyap-criteria}
若存在标量函数 $V(\mathbf{x})$，对所有 $\mathbf{x}$ 有连续一阶偏导、且\emph{正定}（$V(\mathbf{0})=0$、$\mathbf{x}\neq\mathbf{0}$ 时 $V>0$），则按 $\dot V(\mathbf{x})$ 的符号有：
\begin{enumerate}
  \item \textbf{稳定判据}：若 $\dot V$ \emph{半负定}，则 $\mathbf{x}_e$ 李雅普诺夫意义下稳定；
  \item \textbf{渐近稳定判据}：若 $\dot V$ \emph{负定}；或 $\dot V$ 半负定但\emph{对任意 $\mathbf{x}(t_0)\neq\mathbf{0}$，$\dot V$ 沿轨迹不恒为零}，则 $\mathbf{x}_e$ 渐近稳定。若再有 $\|\mathbf{x}\|\to\infty$ 时 $V\to\infty$（\emph{径向无界}），则\emph{大范围渐近稳定}；
  \item \textbf{不稳定判据}：若 $\dot V$ \emph{正定}，则 $\mathbf{x}_e$ 不稳定。
\end{enumerate}
这些条件是\textbf{充分}而非充要的：找到合适的 $V$ 即可下肯定结论；但\emph{找不到}并不能反过来断定不稳定。
\end{theorem}

\begin{insight}[“$\dot V$ 半负定但不恒为零”这一附加条件在补什么]
当 $\dot V$ 只半负定时，$\mathbf{x}\neq\mathbf{0}$ 处可能出现 $\dot V=0$，此时轨迹有两种可能（\cref{fig:vdot-zero}）：(a) $\dot V\equiv0$ 沿轨迹\emph{恒为零}——轨迹被锁在某个曲面 $V(\mathbf{x})=C$ 上，对应非线性系统的极限环或线性系统的临界稳定，\emph{不收敛}；(b) $\dot V$ 只在某瞬时为零（轨迹与某等-$V$ 面\emph{相切}），切过后继续向原点收敛——仍属\emph{渐近}稳定。附加条件“$\dot V$ 不恒为零”正是为了排除情形 (a)、保留情形 (b)。这与导论的 LaSalle 不变原理（\cref{thm:ctrl-lasalle}）说的是同一件事——“最大不变集只含原点”即“$\dot V$ 不恒为零”的几何版本。
\end{insight}

\begin{figure}[H]
\centering
\begin{tikzpicture}[>=Stealth, font=\small]
% (a) trajectory trapped on a level set (limit cycle) -> dotV ≡ 0
\begin{scope}
  \draw[gray!55!black, dashed] (0,0) circle (0.6);
  \draw[gray!55!black, dashed] (0,0) circle (1.45);
  \node[gray!55!black, font=\scriptsize] at (1.15,1.2) {$V=C$};
  \fill (0,0) circle (1.0pt);
  \draw[main, very thick, ->] (1.05,0) arc (0:300:1.05);
  \node[align=center, font=\scriptsize] at (0,-1.95) {(a) $\dot V\equiv0$\\轨迹锁在 $V=C$ 上\\（极限环/临界稳定）};
\end{scope}
% (b) trajectory tangent then spirals in -> still asymptotic
\begin{scope}[xshift=5.0cm]
  \draw[gray!55!black, dashed] (0,0) circle (0.6);
  \draw[gray!55!black, dashed] (0,0) circle (1.45);
  \node[gray!55!black, font=\scriptsize] at (1.15,1.2) {$V=C$};
  \fill (0,0) circle (1.0pt);
  \draw[third, very thick, ->] (1.45,0.1)
        .. controls (1.0,1.0) and (-0.3,1.25) .. (-1.0,0.55)
        .. controls (-1.45,0.0) and (-0.9,-0.7) .. (-0.1,-0.62)
        .. controls (0.5,-0.55) and (0.55,-0.05) .. (0.2,0.2)
        .. controls (0.05,0.32) and (-0.04,0.1) .. (0,0);
  \node[align=center, font=\scriptsize] at (0,-1.95) {(b) $\dot V$ 不恒为零\\与 $V=C$ 相切后继续收敛\\（渐近稳定）};
\end{scope}
\end{tikzpicture}
\caption{$\dot V(\mathbf{x})=0$ 的两种运动（仿刘豹\cite{liubao2006modern} 图 4.5 重绘），解释 \cref{thm:lyap-criteria} 渐近稳定判据里“$\dot V$ 不恒为零”这一附加条件的必要性。}
\label{fig:vdot-zero}
\end{figure}

\subsection{一题多解：在同一系统上打磨“造 \texorpdfstring{$V$}{V}”的直觉}
\label{sec:lyap-examples}

下面几个例题密集训练“先猜 $V$、再算 $\dot V$、按判据下结论”的套路，并贯穿一个核心观察：\emph{$V$ 不唯一，选得好坏直接影响能判出多强的结论}。这是刘豹\cite{liubao2006modern} 第 4 章的精华，全部 TikZ/重排保留。

\begin{example}[最朴素的距离型 $V$]\label{ex:lyap-dist}
非线性系统 $\dot x_1=x_2-x_1(x_1^2+x_2^2)$、$\dot x_2=-x_1-x_2(x_1^2+x_2^2)$，原点是唯一平衡点。取
\[
  V(\mathbf{x})=x_1^2+x_2^2\succ0,
\]
沿轨迹求导：$\dot V=2x_1\dot x_1+2x_2\dot x_2=-2(x_1^2+x_2^2)^2$，负定。且 $\|\mathbf{x}\|\to\infty$ 时 $V\to\infty$，故原点\textbf{大范围渐近稳定}。
\end{example}

\begin{insight}[把 $V$ 读成距离、$\dot V$ 读成趋近速度]
在 \cref{ex:lyap-dist} 中 $V=x_1^2+x_2^2=C$ 是一族以原点为心、半径 $\sqrt C$ 的圆（\cref{fig:level-circles}），代表系统储能；储能越多圆越大、状态离原点越远。$\dot V<0$ 表示状态沿轨线\emph{从外圈穿向内圈}，能量随时间衰减并收敛到原点。\textbf{一句话}：若 $V$ 表示状态到原点的距离，则 $\dot V$ 就是状态\emph{沿轨线趋向原点的速度}。这是理解第二法的最锋利的一把刀。
\end{insight}

\begin{figure}[H]
\centering
\begin{tikzpicture}[>=Stealth, font=\small, scale=1.0]
\draw[->, gray!60!black] (-2.2,0)--(2.2,0) node[right]{$x_1$};
\draw[->, gray!60!black] (0,-2.2)--(0,2.2) node[above]{$x_2$};
\foreach \r in {1.8, 1.2, 0.6}{
  \draw[main!70!black, thick] (0,0) circle (\r);
}
\node[main!60!black, font=\scriptsize] at (1.45,1.45) {$V=C_3$};
\node[main!60!black, font=\scriptsize] at (0.95,0.95) {$C_2$};
\node[main!60!black, font=\scriptsize] at (0.5,0.5) {$C_1$};
% inward spiral trajectory crossing level sets
\draw[third, very thick, ->] (1.95,0.3)
      .. controls (1.2,1.6) and (-1.0,1.7) .. (-1.6,0.4)
      .. controls (-1.9,-0.9) and (0.0,-1.5) .. (0.9,-0.55)
      .. controls (1.35,0.05) and (0.7,0.75) .. (0.0,0.45)
      .. controls (-0.35,0.25) and (-0.25,-0.2) .. (0,0);
\fill (0,0) circle (1.1pt);
\node[third!60!black, font=\scriptsize] at (1.9,0.65) {轨线};
\end{tikzpicture}
\caption{等-$V$ 圆 $V=x_1^2+x_2^2=C_k$（$C_1<C_2<C_3$）与穿越它们向内收敛的轨线（仿刘豹\cite{liubao2006modern} 图 4.6 重绘）。$\dot V<0$ 即轨线从外圈穿向内圈，对应能量耗散、距离递减。}
\label{fig:level-circles}
\end{figure}

\begin{example}[半负定 $\dot V$：从“稳定”补到“渐近稳定”]\label{ex:lyap-semineg}
系统 $\dot{\mathbf{x}}=\big[\begin{smallmatrix}0&1\\-1&-1\end{smallmatrix}\big]\mathbf{x}$，原点唯一平衡。取 $V=x_1^2+x_2^2\succ0$，则
\[
  \dot V=2x_1\dot x_1+2x_2\dot x_2=-2x_2^2,
\]
仅\emph{半}负定（$x_2=0,\ x_1\neq0$ 时 $\dot V=0$）。先得“稳定”。能否升级到渐近稳定？只需验证 $\dot V$ 不沿轨迹恒为零：若 $\dot V=-2x_2^2\equiv0$ 则 $x_2\equiv0$，从而 $\dot x_2\equiv0$；但状态方程 $\dot x_2=-x_1-x_2$ 要求此时 $x_1\equiv0$。故仅在原点才有 $\dot V\equiv0$，沿任一非平凡轨迹 $\dot V$ 不恒为零。配径向无界，原点\textbf{大范围渐近稳定}。这对应 \cref{fig:vdot-zero}b 的“相切后继续收敛”。
\end{example}

\begin{example}[换一个 $V$：直接拿到负定 $\dot V$]\label{ex:lyap-multiV}
仍对 \cref{ex:lyap-semineg} 的系统，另选
\[
  V(\mathbf{x})=\tfrac12\big[(x_1+x_2)^2+2x_1^2+x_2^2\big]
  =(x_1,x_2)\begin{bmatrix}\tfrac32&\tfrac12\\[2pt]\tfrac12&1\end{bmatrix}\!\binom{x_1}{x_2}\succ0,
\]
则 $\dot V=(x_1+x_2)(\dot x_1+\dot x_2)+2x_1\dot x_1+x_2\dot x_2=-(x_1^2+x_2^2)$，\emph{直接负定}，一步判出大范围渐近稳定。\textbf{对比 \cref{ex:lyap-semineg}}：同一系统，“好的” $V$ 让 $\dot V$ 一步负定、省去“不恒为零”的额外论证。这就是\emph{$V$ 选得好不好}的现实意义——它预示了第 5 节“怎样系统地造出好 $V$”的价值。
\end{example}

\begin{example}[临界稳定：闭轨与守恒]\label{ex:lyap-critical}
无阻尼振子 $\dot x_1=x_2$、$\dot x_2=-x_1$（结构不稳定系统的齐次部分）。取 $V=x_1^2+x_2^2\succ0$，则
\[
  \dot V=2x_1\dot x_1+2x_2\dot x_2=2(x_1x_2-x_1x_2)\equiv0.
\]
$\dot V\equiv0$ 表示 $V$ 沿任意轨迹\emph{恒为常数} $V=x_1^2+x_2^2=C$——相轨迹是一族以原点为心、$\sqrt C$ 为半径的圆，等幅振荡。系统李雅普诺夫意义下稳定，但\emph{非}渐近稳定（经典理论称临界稳定，对应 \cref{fig:vdot-zero}a 与 \cref{fig:lyap-ball-domain}a）。要让它渐近稳定，必须\emph{改变系统结构}（加阻尼/反馈）。
\end{example}

\begin{example}[不稳定判据 + 局限性]\label{ex:lyap-limitcycle}
系统 $\dot x_1=x_2$、$\dot x_2=-(1-|x_1|)x_2-x_1$，原点唯一平衡。取 $V=x_1^2+x_2^2$，则 $\dot V=-2x_2^2(1-|x_1|)$。在单位圆内（$|x_1|<1$）$\dot V$ 负定，原点\emph{局部渐近稳定}；在 $|x_1|>1$ 处 $\dot V>0$，系统在单位圆外不稳定。这条 $x_1^2+x_2^2=1$ 上的边界是\emph{不稳定极限环}。\textbf{教训}：第二法的结论是\emph{局部}的——这里同一个 $V$ 在圆内圆外给出相反结论，$V$ 只刻画平衡点某邻域内的运动，对域外一无所知。
\end{example}

\begin{pitfall}[概念误区：找不到 $V$ 就断言“不稳定”]
\cref{thm:lyap-criteria} 是充分条件。\emph{找到}满足判据的 $V$ $\Rightarrow$ 可下肯定结论（稳定/渐近/不稳定）；但\emph{没找到}合适的 $V$，\textbf{不能}据此说系统不稳定——也许只是你没构造出来。这与“证明题没做出来不代表命题为假”同理。正因如此，“怎样系统地造 $V$”（第 5 节）与“存在性逆定理”（交 \cref{ch:lyapunov-stability}）才如此重要。
\end{pitfall}

% ════════════════════════════════════════════════════════════════
\section{李雅普诺夫方法在线性系统中的应用}
\label{sec:lyap-linear-app}

线性系统是李雅普诺夫理论最干净的特例，也是后续\emph{极点配置闭环、观测器误差、LQR 闭环}稳定性分析的工具。本节把导论\cref{thm:ctrl-lyap-eq} 只给结论的线性李氏方程\emph{亲手推一遍、算几个例子}，并补上时变与离散版（现控02 互校）。

\subsection{连续定常系统：李雅普诺夫方程的推导与计算}
\label{sec:lyap-eq-cont}

\begin{theorem}[线性定常系统渐近稳定判据]\label{thm:lyap-eq-cont}
线性定常系统 $\dot{\mathbf{x}}=\mathbf{A}\mathbf{x}$ 的平衡点 $\mathbf{x}_e=\mathbf{0}$ \textbf{大范围渐近稳定}的充要条件是：对\emph{任意}给定的正定实对称阵 $\mathbf{Q}\succ0$，存在正定实对称阵 $\mathbf{P}\succ0$ 满足\textbf{李雅普诺夫方程}
\begin{equation}
  \boxed{\ \mathbf{A}^\top\mathbf{P}+\mathbf{P}\mathbf{A}=-\mathbf{Q}\ }
  \label{eq:lyap-eq}
\end{equation}
此时 $V(\mathbf{x})=\mathbf{x}^\top\mathbf{P}\mathbf{x}$ 是系统的李雅普诺夫函数。这一条件与“$\mathbf{A}$ 特征值均具负实部”\emph{等价}。
\end{theorem}

\begin{proof}[必要性：$\mathbf{A}$ 稳定 $\Rightarrow$ 解存在且正定]
设 $\mathbf{Q}\succ0$，构造
\begin{equation}
  \mathbf{P}=\int_0^{+\infty}e^{\mathbf{A}^\top t}\,\mathbf{Q}\,e^{\mathbf{A}t}\,\mathrm dt.
  \label{eq:lyap-P-int}
\end{equation}
因 $\mathbf{A}$ 的特征值实部均负，$\|e^{\mathbf{A}t}\|$ 指数衰减，积分收敛；被积矩阵正定故 $\mathbf{P}\succ0$。代入验证（用 $\frac{\mathrm d}{\mathrm dt}e^{\mathbf{A}t}=\mathbf{A}e^{\mathbf{A}t}$）：
\begin{align}
  \mathbf{A}^\top\mathbf{P}+\mathbf{P}\mathbf{A}
  &=\int_0^{+\infty}\!\Big(\mathbf{A}^\top e^{\mathbf{A}^\top t}\mathbf{Q}e^{\mathbf{A}t}+e^{\mathbf{A}^\top t}\mathbf{Q}e^{\mathbf{A}t}\mathbf{A}\Big)\mathrm dt
   =\int_0^{+\infty}\frac{\mathrm d}{\mathrm dt}\Big(e^{\mathbf{A}^\top t}\mathbf{Q}e^{\mathbf{A}t}\Big)\mathrm dt\notag\\
  &=\Big[e^{\mathbf{A}^\top t}\mathbf{Q}e^{\mathbf{A}t}\Big]_0^{+\infty}=\mathbf{0}-\mathbf{Q}=-\mathbf{Q},
\end{align}
其中上限处 $e^{\mathbf{A}t}\to\mathbf{0}$。
\end{proof}

\begin{proof}[充分性：解正定 $\Rightarrow$ $\mathbf{A}$ 稳定]
在复数域上讨论（$\mathbf{A}$ 的特征值可能复）。在 $\mathbb{C}^n$ 中定义内积 $\langle\mathbf{x},\mathbf{y}\rangle=\mathbf{x}^\top\mathbf{P}\bar{\mathbf{y}}$。设 $\lambda\in\sigma(\mathbf{A})$、$\mathbf{x}\neq\mathbf{0}$ 为对应特征向量（$\mathbf{A}\mathbf{x}=\lambda\mathbf{x}$）。一方面
\[
  \langle\mathbf{A}\mathbf{x},\mathbf{x}\rangle+\langle\mathbf{x},\mathbf{A}\mathbf{x}\rangle
  =\mathbf{x}^\top(\mathbf{A}^\top\mathbf{P}+\mathbf{P}\mathbf{A})\bar{\mathbf{x}}=-\mathbf{x}^\top\mathbf{Q}\bar{\mathbf{x}}<0;
\]
另一方面
\[
  \langle\mathbf{A}\mathbf{x},\mathbf{x}\rangle+\langle\mathbf{x},\mathbf{A}\mathbf{x}\rangle
  =(\lambda+\bar\lambda)\,\mathbf{x}^\top\mathbf{P}\bar{\mathbf{x}}=2\,\mathrm{Re}(\lambda)\,\mathbf{x}^\top\mathbf{P}\bar{\mathbf{x}}.
\]
因 $\mathbf{x}^\top\mathbf{P}\bar{\mathbf{x}}>0$，比较两式得 $\mathrm{Re}(\lambda)<0$，即所有特征值落在左半平面，$\mathbf{A}$ 稳定。
\end{proof}

\noindent\textbf{$\dot V$ 的代数核.}\;取 $V=\mathbf{x}^\top\mathbf{P}\mathbf{x}$，沿 $\dot{\mathbf{x}}=\mathbf{A}\mathbf{x}$ 求导：
\begin{equation}
  \dot V=\dot{\mathbf{x}}^\top\mathbf{P}\mathbf{x}+\mathbf{x}^\top\mathbf{P}\dot{\mathbf{x}}
  =\mathbf{x}^\top(\mathbf{A}^\top\mathbf{P}+\mathbf{P}\mathbf{A})\mathbf{x}=-\mathbf{x}^\top\mathbf{Q}\mathbf{x},
  \label{eq:Vdot-core}
\end{equation}
要 $\dot V$ 负定即要 $\mathbf{Q}\succ0$。\emph{实用流程}是反着走：先指定 $\mathbf{Q}$（常取 $\mathbf{Q}=\mathbf{I}$），解 \cref{eq:lyap-eq} 得 $\mathbf{P}$，再用希尔维斯特判据验 $\mathbf{P}\succ0$、从而判稳定。之所以“先 $\mathbf{Q}$ 后 $\mathbf{P}$”而非反向，是因为给定 $\mathbf{Q}$ 解 $\mathbf{P}$ 是\emph{线性}方程组、唯一可解；反向先取 $\mathbf{P}$ 再算 $\mathbf{Q}$ 虽简单却难保正定。

\begin{pitfall}[证明陷阱：误以为“$\mathbf{A}$ 稳定 $\Rightarrow$ $\mathbf{A}+\mathbf{A}^\top\prec0$”]
$\mathbf{A}^\top\mathbf{P}+\mathbf{P}\mathbf{A}$ 是 $\dot V$ 的核（\cref{eq:Vdot-core}），它\emph{对称}（即便 $\mathbf{A}$ 不对称）。新手常误以为“稳定就有 $\mathbf{A}+\mathbf{A}^\top\prec0$”——\textbf{错}：$\mathbf{A}+\mathbf{A}^\top$ 只是 $\mathbf{P}=\mathbf{I}$ 这一\emph{特例}的核，稳定矩阵未必满足（如 $\mathbf{A}=\big[\begin{smallmatrix}-1&10\\0&-1\end{smallmatrix}\big]$ 稳定但 $\mathbf{A}+\mathbf{A}^\top$ 不负定）。正确做法是\emph{对给定 $\mathbf{Q}$ 解出 $\mathbf{P}$}（\cref{eq:lyap-P-int}），不能直接取 $\mathbf{P}=\mathbf{I}$。仅当通过相似变换把 $\mathbf{A}$ 对角化为特征值阵 $\boldsymbol\Lambda$、取 $V=\|\mathbf{x}\|^2$ 时才有 $\mathbf{Q}=-(\boldsymbol\Lambda^\top+\boldsymbol\Lambda)=-2\boldsymbol\Lambda$，此时“$\boldsymbol\Lambda$ 对角元（特征值）全负 $\iff$ $\mathbf{Q}\succ0$”才成立。% OCR: 刘豹源此处"设 A=A、Q=-2A"系 \Lambda 误识为 A，已据勘误改为对角阵 \Lambda。
（导论 \cref{thm:ctrl-lyap-eq} 的对应陷阱栏与此一致。）
\end{pitfall}

\begin{example}[$\mathbf{Q}=\mathbf{I}$ 解 $\mathbf{P}$ 判稳定]\label{ex:lyap-solve-P}
系统 $\dot{\mathbf{x}}=\big[\begin{smallmatrix}0&1\\-2&-3\end{smallmatrix}\big]\mathbf{x}$。设 $\mathbf{P}=\big[\begin{smallmatrix}p_{11}&p_{12}\\p_{12}&p_{22}\end{smallmatrix}\big]$、$\mathbf{Q}=\mathbf{I}$，代入 $\mathbf{A}^\top\mathbf{P}+\mathbf{P}\mathbf{A}=-\mathbf{I}$ 展开、令对应元相等，解得
\[
  \mathbf{P}=\begin{bmatrix}\tfrac54&\tfrac14\\[2pt]\tfrac14&\tfrac14\end{bmatrix}.
\]
希尔维斯特判据：$\Delta_1=\tfrac54>0$、$\Delta_2=\det\mathbf{P}=\tfrac14>0$，故 $\mathbf{P}\succ0$，系统\textbf{大范围渐近稳定}。亦可直接看 $V=\tfrac14(5x_1^2+2x_1x_2+x_2^2)\succ0$、$\dot V=-(x_1^2+x_2^2)\prec0$ 得同样结论。
\end{example}

\begin{example}[用李氏方程定参数范围：半正定 $\mathbf{Q}$ 的妙用]\label{ex:lyap-param}
系统 $\dot{\mathbf{x}}=\big[\begin{smallmatrix}0&1&0\\0&-2&1\\-K&0&-1\end{smallmatrix}\big]\mathbf{x}$，求增益 $K$ 的稳定范围。因 $\det\mathbf{A}\neq0$，原点唯一平衡。这里巧选\emph{半}正定
\[
  \mathbf{Q}=\begin{bmatrix}0&0&0\\0&0&0\\0&0&1\end{bmatrix},\qquad
  \dot V=-\mathbf{x}^\top\mathbf{Q}\mathbf{x}=-x_3^2.
\]
% OCR: 刘豹源 \dot V=x_3^2 漏负号，且把 \dot V 误印为 V；据勘误应为 \dot V=-x_3^2。
半正定 $\mathbf{Q}$ 合法的前提是 $\dot V$ 沿轨迹不恒为零：$\dot V\equiv0$ 要 $x_3\equiv0$，由状态方程逐级推得 $x_1\equiv x_2\equiv0$，即仅原点处恒零——故 $\mathbf{Q}$ 可半正定。解 $\mathbf{A}^\top\mathbf{P}+\mathbf{P}\mathbf{A}=-\mathbf{Q}$ 得
\[
  \mathbf{P}=\frac{1}{12-2K}\begin{bmatrix}K^2+12K&6K&0\\6K&3K&K\\0&K&6\end{bmatrix},
\]
其正定的充要条件是 $12-2K>0$ 且 $K>0$，即
\[
  0<K<6.
\]
故 $0<K<6$ 时原点大范围渐近稳定。\textbf{方法论意义}：李氏方程不只判“稳不稳”，还能反解出\emph{保稳定的参数区间}——这是它比单纯算特征值更有用的地方，也预告了它在控制器/观测器增益设计（\cref{ch:linear-synthesis,ch:optimal-control-basics}）中的角色。
\end{example}

\subsection{时变与离散系统（互校：现控02）}
\label{sec:lyap-tv-disc}

李雅普诺夫第二法的力量在于它\emph{无缝}推广到时变与离散系统，给出统一形式的判据。这部分刘豹与现控02 表述一致，互为第二信源。

\begin{theorem}[线性时变连续系统渐近稳定判据]\label{thm:lyap-tv}
线性时变系统 $\dot{\mathbf{x}}=\mathbf{A}(t)\mathbf{x}$ 在 $\mathbf{x}_e=\mathbf{0}$ 处大范围渐近稳定的充要条件是：对任意给定连续对称正定 $\mathbf{Q}(t)$，存在连续对称正定 $\mathbf{P}(t)$ 满足\emph{矩阵微分方程}
\begin{equation}
  \dot{\mathbf{P}}(t)=-\mathbf{A}^\top(t)\mathbf{P}(t)-\mathbf{P}(t)\mathbf{A}(t)-\mathbf{Q}(t),
  \label{eq:lyap-tv}
\end{equation}
且李氏函数为 $V(\mathbf{x},t)=\mathbf{x}^\top\mathbf{P}(t)\mathbf{x}$。
\end{theorem}
\begin{proof}[要点]
取 $V=\mathbf{x}^\top\mathbf{P}(t)\mathbf{x}$，全导数 $\dot V=\mathbf{x}^\top\big[\mathbf{A}^\top(t)\mathbf{P}(t)+\dot{\mathbf{P}}(t)+\mathbf{P}(t)\mathbf{A}(t)\big]\mathbf{x}=-\mathbf{x}^\top\mathbf{Q}(t)\mathbf{x}$，要 $\dot V$ 负定即要 $\mathbf{Q}(t)\succ0$。\cref{eq:lyap-tv} 是黎卡提（Riccati）矩阵微分方程的特例——这一点为 \cref{ch:lqr-lqg-riccati} 的有限时域 LQR 埋下伏笔（彼处的代价权重矩阵满足同型方程）。
\end{proof}

\begin{note}
\textbf{记号双关提示（控制 vs 估计）.}\;\cref{eq:lyap-tv} 与导论的连续 Riccati 微分方程（\cref{sec:ctrl-lqr-rde}）同型，而后者又与卡尔曼滤波协方差传播（\cref{ch:eskf}）形式对偶。这里 $\mathbf{P}(t)$ 是李氏/代价权重，估计侧的协方差另记 $\boldsymbol\Sigma(t)$——\emph{形式相同、物理含义对偶}（控制是“代价 cost-to-go”，估计是“不确定度”），切勿因共用 Riccati 形式而混淆二者语义。
\end{note}

\begin{theorem}[线性定常离散系统渐近稳定判据]\label{thm:lyap-disc}
线性定常离散系统 $\mathbf{x}(k+1)=\mathbf{G}\mathbf{x}(k)$ 在 $\mathbf{x}_e=\mathbf{0}$ 处渐近稳定的充要条件是 $\mathbf{G}$ 的特征根均在单位开圆盘内（谱半径 $<1$）；等价地，对任意 $\mathbf{Q}\succ0$，存在 $\mathbf{P}\succ0$ 满足\textbf{离散李雅普诺夫方程}
\begin{equation}
  \mathbf{G}^\top\mathbf{P}\mathbf{G}-\mathbf{P}=-\mathbf{Q}.
  \label{eq:lyap-disc}
\end{equation}
此时 $V[\mathbf{x}(k)]=\mathbf{x}^\top(k)\mathbf{P}\mathbf{x}(k)$ 为李氏函数，判稳定看的是\emph{一步差分}
\begin{equation}
  \Delta V=V[\mathbf{x}(k+1)]-V[\mathbf{x}(k)]=\mathbf{x}^\top(k)\big(\mathbf{G}^\top\mathbf{P}\mathbf{G}-\mathbf{P}\big)\mathbf{x}(k)=-\mathbf{x}^\top(k)\mathbf{Q}\mathbf{x}(k).
  \label{eq:lyap-disc-delta}
\end{equation}
\end{theorem}

\begin{insight}[连续 $\dot V$ $\leftrightarrow$ 离散 $\Delta V$ 的平行]
离散版把连续的“导数 $\dot V$ 负定”整体替换为“一步差分 $\Delta V$ 负定”，把李氏方程的 $\mathbf{A}^\top\mathbf{P}+\mathbf{P}\mathbf{A}$ 换成 $\mathbf{G}^\top\mathbf{P}\mathbf{G}-\mathbf{P}$，把稳定区从“左半平面”换到“单位圆内”。这套\emph{连续—离散对照}贯穿全书：MPC 的终端代价（\cref{sec:ctrl-mpc}）用的就是离散李氏函数。其充分性证明（对 \cref{eq:lyap-disc} 逐次左右乘 $\mathbf{G}$ 求和，利用 $\lim_{n\to\infty}(\mathbf{G}^{n})^\top\mathbf{P}\mathbf{G}^{n}=\mathbf{0}$ 得 $\mathbf{P}=\sum_{k=1}^\infty(\mathbf{G}^k)^\top\mathbf{Q}\mathbf{G}^k\succ0$）与连续情形完全平行，留作练习。
\end{insight}

\begin{example}[对角离散系统：极点落单位圆内]\label{ex:lyap-disc-diag}
$\mathbf{x}(k+1)=\big[\begin{smallmatrix}\lambda_1&0\\0&\lambda_2\end{smallmatrix}\big]\mathbf{x}(k)$。取 $\mathbf{Q}=\mathbf{I}$，由 \cref{eq:lyap-disc} 解得 $\mathbf{P}=\mathrm{diag}\big(\tfrac{1}{1-\lambda_1^2},\tfrac{1}{1-\lambda_2^2}\big)$。$\mathbf{P}\succ0$ 的充要条件是 $|\lambda_1|<1$ 且 $|\lambda_2|<1$——极点落在单位圆内，与采样控制系统的稳定判据完全一致。
\end{example}

\begin{theorem}[线性时变离散系统渐近稳定判据]\label{thm:lyap-tv-disc}
线性时变离散系统 $\mathbf{x}(k+1)=\mathbf{G}(k)\mathbf{x}(k)$ 在 $\mathbf{x}_e=\mathbf{0}$ 处大范围渐近稳定的充要条件是：对任意给定的对称正定矩阵序列 $\{\mathbf{Q}(k)\succ0\}$，存在对称正定矩阵序列 $\{\mathbf{P}(k)\succ0\}$ 满足\emph{差分型}离散李雅普诺夫方程
\begin{equation}
  \mathbf{G}^\top(k)\,\mathbf{P}(k+1)\,\mathbf{G}(k)-\mathbf{P}(k)=-\mathbf{Q}(k),
  \label{eq:lyap-tv-disc}
\end{equation}
此时李氏函数取\emph{时变}二次型 $V[\mathbf{x}(k),k]=\mathbf{x}^\top(k)\mathbf{P}(k)\mathbf{x}(k)$。
\end{theorem}
\begin{proof}[要点]
沿轨迹取一步差分 $\Delta V=V[\mathbf{x}(k{+}1),k{+}1]-V[\mathbf{x}(k),k]=\mathbf{x}^\top(k{+}1)\mathbf{P}(k{+}1)\mathbf{x}(k{+}1)-\mathbf{x}^\top(k)\mathbf{P}(k)\mathbf{x}(k)$，代入 $\mathbf{x}(k{+}1)=\mathbf{G}(k)\mathbf{x}(k)$ 得 $\Delta V=\mathbf{x}^\top(k)\big[\mathbf{G}^\top(k)\mathbf{P}(k{+}1)\mathbf{G}(k)-\mathbf{P}(k)\big]\mathbf{x}(k)=-\mathbf{x}^\top(k)\mathbf{Q}(k)\mathbf{x}(k)\prec0$。它把 \cref{thm:lyap-disc} 的\emph{定常}离散方程 $\mathbf{G}^\top\mathbf{P}\mathbf{G}-\mathbf{P}=-\mathbf{Q}$ 推广到逐步变化的 $\mathbf{G}(k),\mathbf{P}(k)$，与 \cref{thm:lyap-tv} 连续时变情形的 $\dot{\mathbf{P}}(t)$ 方程恰为离散对偶（导数 $\dot V$ $\to$ 一步差分 $\Delta V$、微分方程 $\to$ 逐步递推方程）。
\end{proof}

\begin{note}
\textbf{四态李氏方程一览（互校：刘豹§4.4）.}\;线性李雅普诺夫方程按“连续/离散 $\times$ 定常/时变”恰成四态，至此已补齐：
\begin{center}\small
\begin{tabular}{lll}
\toprule
 & 定常 & 时变 \\
\midrule
连续 & $\mathbf{A}^\top\mathbf{P}+\mathbf{P}\mathbf{A}=-\mathbf{Q}$（\cref{eq:lyap-eq}） & $\dot{\mathbf{P}}=-\mathbf{A}^\top(t)\mathbf{P}-\mathbf{P}\mathbf{A}(t)-\mathbf{Q}(t)$（\cref{eq:lyap-tv}） \\
离散 & $\mathbf{G}^\top\mathbf{P}\mathbf{G}-\mathbf{P}=-\mathbf{Q}$（\cref{eq:lyap-disc}） & $\mathbf{G}^\top(k)\mathbf{P}(k{+}1)\mathbf{G}(k)-\mathbf{P}(k)=-\mathbf{Q}(k)$（\cref{eq:lyap-tv-disc}） \\
\bottomrule
\end{tabular}
\end{center}
\emph{读法}：连续侧的导数 $\dot V$ 对应离散侧的一步差分 $\Delta V$；定常侧的代数方程对应时变侧的逐步递推方程。时变离散判据在采样数据系统、时变 MPC 的终端代价分析中用到，且与 \cref{ch:eskf} 离散时变协方差传播 $\boldsymbol\Sigma(k{+}1)=\mathbf{G}(k)\boldsymbol\Sigma(k)\mathbf{G}^\top(k)+\cdots$ 形式对偶——这是“估计--控制对偶”在\emph{离散时变}层面的再次体现。
\end{note}

% ════════════════════════════════════════════════════════════════
\section{怎样构造李雅普诺夫函数}
\label{sec:lyap-construct}

\paragraph{这是第二法真正的难关。}
李雅普诺夫稳定性理论\emph{本身}并不提供构造 $V$ 的一般方法——这曾长期阻碍第二法的推广。本节给出两条\emph{可操作}的造 $V$ 流程：克拉索夫斯基（雅可比）法与变量梯度法。它们都是“把线性系统的二次型李氏函数推广到非线性”的尝试。\textbf{这一整节是导论完全没有覆盖的差异化干货}，也是本章作为“做透层”最核心的增量。

\subsection{克拉索夫斯基（雅可比）法}
\label{sec:lyap-krasovskii}

\paragraph{核心思想：不用 $\mathbf{x}$ 而用 $\dot{\mathbf{x}}=\mathbf{f}(\mathbf{x})$ 凑二次型。}
对某些非线性系统，与其用状态 $\mathbf{x}$ 构造李氏函数，不如\emph{用 $\dot{\mathbf{x}}$ 的二次型}。克拉索夫斯基（Krasovskii）法就把 $V$ 选成 $\mathbf{f}(\mathbf{x})$ 的（带权）欧几里得范数平方\cite{slotine1991applied}。

\begin{theorem}[克拉索夫斯基法]\label{thm:krasovskii}
设 $\dot{\mathbf{x}}=\mathbf{f}(\mathbf{x})$、$\mathbf{f}(\mathbf{0})=\mathbf{0}$、$\mathbf{f}$ 对各 $x_i$ 可微，雅可比 $\mathbf{J}(\mathbf{x})=\partial\mathbf{f}/\partial\mathbf{x}$。任给正定实对称阵 $\mathbf{P}$，若矩阵
\begin{equation}
  \mathbf{Q}(\mathbf{x})=-\big[\mathbf{J}^\top(\mathbf{x})\mathbf{P}+\mathbf{P}\mathbf{J}(\mathbf{x})\big]
  \label{eq:kras-Q}
\end{equation}
对所有 $\mathbf{x}\neq\mathbf{0}$ 正定，则原点渐近稳定，且
\begin{equation}
  V(\mathbf{x})=\dot{\mathbf{x}}^\top\mathbf{P}\dot{\mathbf{x}}=\mathbf{f}^\top(\mathbf{x})\,\mathbf{P}\,\mathbf{f}(\mathbf{x})
  \label{eq:kras-V}
\end{equation}
是一个李氏函数。% OCR: 刘豹源 V=x^T P \dot x 系首因子 \dot x 的导数点漏识，据勘误及式(4.51)应为 \dot x^T P \dot x。
若 $\|\mathbf{x}\|\to\infty$ 时 $V\to\infty$ 则大范围渐近稳定。特别取 $\mathbf{P}=\mathbf{I}$ 得\textbf{克拉索夫斯基表达式}
\begin{equation}
  \mathbf{Q}(\mathbf{x})=-\big[\mathbf{J}^\top(\mathbf{x})+\mathbf{J}(\mathbf{x})\big],\qquad V(\mathbf{x})=\mathbf{f}^\top(\mathbf{x})\mathbf{f}(\mathbf{x}).
  \label{eq:kras-I}
\end{equation}
\end{theorem}

\begin{proof}[推导]
取 $V=\mathbf{f}^\top\mathbf{P}\mathbf{f}\succ0$（$\mathbf{P}\succ0$）。因 $\mathbf{f}$ 是 $\mathbf{x}$ 的显函数、不显含 $t$，有链式关系 $\dot{\mathbf{f}}=\dfrac{\partial\mathbf{f}}{\partial\mathbf{x}}\dot{\mathbf{x}}=\mathbf{J}(\mathbf{x})\mathbf{f}(\mathbf{x})$。沿轨迹求导：
\begin{equation}
  \dot V=\dot{\mathbf{f}}^\top\mathbf{P}\mathbf{f}+\mathbf{f}^\top\mathbf{P}\dot{\mathbf{f}}
  =\mathbf{f}^\top\big[\mathbf{J}^\top\mathbf{P}+\mathbf{P}\mathbf{J}\big]\mathbf{f}=-\mathbf{f}^\top\mathbf{Q}(\mathbf{x})\mathbf{f}.
\end{equation}
要 $\dot V$ 负定即要 $\mathbf{Q}(\mathbf{x})\succ0$。
\end{proof}

\begin{note}
\textbf{现控02 的独立表述（互校）.}\;现控02\cite{zhang2017modern} 把同一定理写成赫米特形式：定义 $\hat{\mathbf{F}}(\mathbf{x})=\mathbf{F}^*(\mathbf{x})+\mathbf{F}(\mathbf{x})$（$\mathbf{F}=\mathbf{J}$ 为雅可比，$\mathbf{F}^*$ 为共轭转置），若 $\hat{\mathbf{F}}$ 负定则 $V=\mathbf{f}^*(\mathbf{x})\mathbf{f}(\mathbf{x})$ 为李氏函数；更一般地 $\hat{\mathbf{F}}=\mathbf{F}^*\mathbf{P}+\mathbf{P}\mathbf{F}+\mathbf{Q}$、$V=\mathbf{f}^*\mathbf{P}\mathbf{f}$。这与刘豹的实矩阵表述 \cref{eq:kras-Q,eq:kras-V} 完全等价（实系统下 $\mathbf{F}^*=\mathbf{J}^\top$、$\hat{\mathbf{F}}=-\mathbf{Q}$），互为第二信源。现控02 还点明：该定理对\emph{非线性}系统给大范围渐近稳定的\emph{充分}条件、对线性系统则\emph{退化为充要}（与 \cref{thm:lyap-eq-cont} 一致）。
\end{note}

\begin{insight}[一个结构性的必要条件]
要 $\mathbf{Q}(\mathbf{x})\succ0$，必须 $\mathbf{J}(\mathbf{x})$ 主对角元\emph{不恒为零}。若某 $f_i(\mathbf{x})$ 不含 $x_i$，则 $\partial f_i/\partial x_i\equiv0$，主对角第 $i$ 元恒零，$\mathbf{Q}$ 不可能正定——此时克拉索夫斯基法\emph{用不上}（不代表系统不稳定，只是这个工具失效）。这给了一个\emph{开算前的快速否决}：先扫一眼对角元。
\end{insight}

\begin{example}[克拉索夫斯基法]\label{ex:krasovskii}
系统 $\dot x_1=-3x_1+x_2$、$\dot x_2=x_1-x_2-x_2^3$。雅可比
\[
  \mathbf{J}(\mathbf{x})=\begin{bmatrix}-3&1\\1&-1-3x_2^2\end{bmatrix},\qquad
  -\mathbf{Q}(\mathbf{x})=\mathbf{J}^\top+\mathbf{J}=\begin{bmatrix}-6&2\\2&-2-6x_2^2\end{bmatrix}.
\]
希尔维斯特判 $\mathbf{Q}\succ0$：$\Delta_1=6>0$、$\Delta_2=\big|\begin{smallmatrix}6&-2\\-2&2+6x_2^2\end{smallmatrix}\big|=8+36x_2^2>0$，故对 $\mathbf{x}\neq\mathbf{0}$ 处处正定。又 $V=\mathbf{f}^\top\mathbf{f}=(-3x_1+x_2)^2+(x_1-x_2-x_2^3)^2\to\infty$（$\|\mathbf{x}\|\to\infty$），故原点\textbf{大范围渐近稳定}。（现控02 例 5.4.1 用 $\dot x_1=-x_1$、$\dot x_2=x_1-x_2-x_2^3$ 给出同型结论，可对照。）
\end{example}

\begin{pitfall}[陷阱：克拉索夫斯基的“处处正定 $\mathbf{Q}$”条件过严]
\cref{thm:krasovskii} 要求 $\mathbf{Q}(\mathbf{x})$ 对\emph{所有} $\mathbf{x}\neq\mathbf{0}$ 正定，这个条件相当苛刻——许多本来稳定的非线性系统并不满足它。一旦 $\hat{\mathbf{F}}$（或 $\mathbf{Q}$）某处不负定，克拉索夫斯基法对稳定性\emph{不提供任何信息}（既不说稳也不说不稳）。它只是充分条件，不要把“克拉索夫斯基法失败”误读为“系统不稳定”。
\end{pitfall}

\subsection{变量梯度法}
\label{sec:lyap-vargrad}

\paragraph{核心思想：先猜梯度 $\nabla V$、再线积分回 $V$。}
变量梯度法（Schultz–Gibson 法，1962）是另一条更灵活的造 $V$ 流程\cite{slotine1991applied}。它的出发点是：若理想的 $V(\mathbf{x})$ 存在，则其梯度 $\nabla V=\partial V/\partial\mathbf{x}$ 必存在且唯一，且
\begin{equation}
  \dot V=\sum_{i=1}^n\frac{\partial V}{\partial x_i}\dot x_i=(\nabla V)^\top\dot{\mathbf{x}}.
  \label{eq:vargrad-Vdot}
\end{equation}
于是\emph{倒过来做}：先假设 $\nabla V$ 是一个带待定系数的向量，由 \cref{eq:vargrad-Vdot} 算出 $\dot V$ 并强制其负定来定系数，最后\emph{线积分}还原 $V$：
\begin{equation}
  V(\mathbf{x})=\int_0^{\mathbf{x}}(\nabla V)^\top\,\mathrm d\mathbf{x}.
  \label{eq:vargrad-V}
\end{equation}

\paragraph{无关路径的条件：旋度为零。}
要让线积分 \cref{eq:vargrad-V} 的结果\emph{与积分路径无关}（从而 $V$ 唯一），必须 $\nabla V$ 旋度为零，等价于其雅可比对称：
\begin{equation}
  \frac{\partial(\nabla V)_i}{\partial x_j}=\frac{\partial(\nabla V)_j}{\partial x_i},\qquad i,j=1,\dots,n,
  \label{eq:curl-zero}
\end{equation}
% OCR: 现控02 源旋度方程分母误印 \partial \dot x_1，据勘误应为 \partial x_1（不带导数点）。
即 $n$ 维系统有 $n(n-1)/2$ 个旋度方程。满足后取最简\emph{逐点积分路径}（先沿 $x_1$ 轴、再沿 $x_2$ 轴、$\dots$）：
\begin{equation}
  V(\mathbf{x})=\int_0^{x_1}\!(\nabla V)_1\big|_{x_2=\cdots=x_n=0}\mathrm dx_1
  +\int_0^{x_2}\!(\nabla V)_2\big|_{x_1\text{ fixed},\,x_3=\cdots=0}\mathrm dx_2+\cdots.
  \label{eq:vargrad-path}
\end{equation}

\begin{algo}[变量梯度法步骤]\label{algo:vargrad}
\begin{enumerate}
  \item \textbf{设梯度}：令 $(\nabla V)_i=\sum_j a_{ij}x_j$，系数 $a_{ij}$ 可为常数 / $t$ 的函数 / 状态的函数；常把 $a_{nn}$ 取常数、部分 $a_{ij}$ 取零；
  \item \textbf{算 $\dot V$}：由 \cref{eq:vargrad-Vdot} 得 $\dot V=(\nabla V)^\top\dot{\mathbf{x}}$；
  \item \textbf{定系数}：按“$\dot V$ 负定（或半负定）”\emph{且}满足 $n(n-1)/2$ 个旋度方程 \cref{eq:curl-zero}，确定余下系数；
  \item \textbf{线积分}：由 \cref{eq:vargrad-path} 还原 $V(\mathbf{x})$，校验其正定性；
  \item \textbf{判范围}：若 $\|\mathbf{x}\|\to\infty$ 时 $V\to\infty$ 则大范围渐近稳定，否则定出渐近稳定区域。
\end{enumerate}
\end{algo}

\begin{example}[变量梯度法：一题两解，比较稳定区域]\label{ex:vargrad}
对系统 $\dot x_1=-x_1$、$\dot x_2=-x_2+x_1x_2^2$ 造 $V$。设 $\nabla V=\binom{a_{11}x_1+a_{12}x_2}{a_{21}x_1+a_{22}x_2}$。

\textbf{解一（简单系数）.}\;试取 $a_{11}=a_{22}=1$、$a_{12}=a_{21}=0$，则
\[
  \dot V=(\nabla V)^\top\dot{\mathbf{x}}=-x_1^2-(1-x_1x_2)x_2^2,
\]
当 $x_1x_2<1$ 时 $\dot V$ 负定。旋度方程 $\partial(\nabla V)_1/\partial x_2=\partial(\nabla V)_2/\partial x_1=0$ 自动满足。线积分得 $V=\tfrac12(x_1^2+x_2^2)\succ0$。故在 $x_1x_2<1$ 范围内原点渐近稳定。

\textbf{解二（含状态的系数）.}\;改取 $a_{11}=1$、$a_{12}=x_2^2$、$a_{21}=3x_2^2$、$a_{22}=3$，则 $\nabla V=\binom{x_1+x_2^3}{3x_1x_2^2+3x_2}$，旋度方程 $\partial(\nabla V)_1/\partial x_2=\partial(\nabla V)_2/\partial x_1=3x_2^2$ 满足。$\dot V$ 负定要 $0<x_1x_2<\tfrac13$，线积分得 $V=\tfrac12x_1^2+\tfrac32x_2^2+x_1x_2^3$，在该范围内正定。

\textbf{比较}：两个 $V$ 不同，给出的渐近稳定区域也不同——解一的 $x_1x_2<1$ 比解二的 $0<x_1x_2<\tfrac13$ \emph{宽得多}（\cref{fig:stab-region}）。这再次印证 \cref{ex:lyap-multiV} 的教训：\emph{$V$ 的选择直接决定结论的强弱}，没有“唯一正确”的 $V$，但有“更好”的 $V$。
\end{example}

\begin{figure}[H]
\centering
\begin{tikzpicture}[>=Stealth, font=\small, scale=1.0]
\draw[->, gray!60!black] (-2.6,0)--(2.6,0) node[right]{$x_1$};
\draw[->, gray!60!black] (0,-2.4)--(0,2.4) node[above]{$x_2$};
% region x1 x2 < 1 : everything except two corners x1x2>=1 (first & third quadrants beyond hyperbola)
% hyperbola x1 x2 = 1 (first quadrant) and =1 (third)
\draw[main, thick, domain=0.45:2.4, samples=60] plot (\x, {1/\x});
\draw[main, thick, domain=-2.4:-0.45, samples=60] plot (\x, {1/\x});
\node[main!70!black, font=\scriptsize] at (2.0,0.75) {$x_1x_2=1$};
% hyperbola x1 x2 = 1/3
\draw[third, thick, domain=0.18:2.4, samples=70] plot (\x, {0.3333/\x});
\node[third!70!black, font=\scriptsize] at (1.85,0.05) {$x_1x_2=\tfrac13$};
% shade region 0<x1x2<1/3 in first quadrant (between axes and 1/3 hyperbola)
\fill[third!18, domain=0.18:2.4, samples=70]
   plot (\x,{0.3333/\x}) -- (2.4,0.001) -- (0.18,0.001) -- cycle;
\node[align=center, font=\scriptsize, main!60!black] at (-1.4,1.55) {解一区域\\$x_1x_2<1$\\（更宽）};
\node[align=center, font=\scriptsize, third!55!black] at (1.05,0.55) {解二\\$0<x_1x_2<\tfrac13$};
\end{tikzpicture}
\caption{\cref{ex:vargrad} 两个李氏函数给出的渐近稳定区域对比（仿刘豹\cite{liubao2006modern} 图 4.9 重绘）。解一的 $x_1x_2<1$（去掉一、三象限远端两块）显著宽于解二阴影区 $0<x_1x_2<\tfrac13$。说明“造 $V$”是有优劣之分的技艺。}
\label{fig:stab-region}
\end{figure}

\begin{example}[变量梯度法用于时变系统]\label{ex:vargrad-tv}
时变系统 $\dot{\mathbf{x}}=\big[\begin{smallmatrix}0&1\\-\frac{1}{t+1}&-10\end{smallmatrix}\big]\mathbf{x}$（$t\ge0$）。取 $\nabla V=\binom{a_{11}x_1}{a_{22}x_2}$（$a_{12}=a_{21}=0$ 满足旋度方程），令 $a_{11}=1$、$a_{22}=t+1$，得 $\nabla V=\binom{x_1}{(t+1)x_2}$，线积分得 $V=\tfrac12[x_1^2+(t+1)x_2^2]\succ0$。其导数 $\dot V=-(10t+10)x_2^2$ 半负定但不恒为零，故原点\textbf{大范围渐近稳定}。这显示变量梯度法对\emph{显含 $t$} 的系数也能处理——把 $a_{ij}$ 取成 $t$ 的函数即可。
\end{example}

\begin{note}
\textbf{对 $V$ 的几点共识（两书一致）.}\;造 $V$ 时记住：(i) $V$ 必正定、对 $\mathbf{x}$ 有连续一阶偏导；(ii) $V$ \emph{非唯一}，但不同 $V$ 不影响“是否稳定”的结论一致性，只影响判出的\emph{范围}；(iii) 最简形式是二次型 $\mathbf{x}^\top\mathbf{P}\mathbf{x}$，但未必都是二次型；(iv) $V$ 只刻画平衡点某邻域内的局部运动，对域外无信息；(v) 造 $V$ 需技巧，故第二法主要用于\emph{其它方法失效或难判}的高阶非线性/时变系统。现控02 的“几点说明”与刘豹的“对 $V(\mathbf{x})$ 的讨论”在这五点上完全吻合（互校）。
\end{note}

% ════════════════════════════════════════════════════════════════
\section{李雅普诺夫函数的进阶用途（现控02）}
\label{sec:lyap-advanced-use}

李雅普诺夫第二法不止于“判稳定”。本节补两项导论未涉及、由现控02\cite{zhang2017modern} 系统给出的用途——它们把 $V$ 从“稳定性裁判”升级为“性能分析与设计工具”，正好衔接最优控制基础（\cref{ch:linear-synthesis,ch:optimal-control-basics}）。

\subsection{用 \texorpdfstring{$V$}{V} 求解参数最优化问题}
\label{sec:lyap-paramopt}

\paragraph{思想：把二次型代价与李氏函数挂钩。}
设系统 $\dot{\mathbf{x}}=\mathbf{A}(\boldsymbol\xi)\mathbf{x}$，$\mathbf{x}(t_0)=\mathbf{x}_0$，$\mathbf{A}(\boldsymbol\xi)$ 含可调参数 $\boldsymbol\xi$。希望在系统渐近稳定（任意初态趋于原点）前提下，使二次型性能指标
\begin{equation}
  J=\int_{t_0}^{+\infty}\mathbf{x}^\top\mathbf{Q}\mathbf{x}\,\mathrm dt
  \label{eq:paramopt-J}
\end{equation}
取极小，从而确定 $\boldsymbol\xi$。\textbf{关键一步}：若 $\mathbf{A}(\boldsymbol\xi)$ 稳定，由 \cref{thm:lyap-eq-cont} 存在 $\mathbf{P}\succ0$ 使 $\mathbf{A}^\top\mathbf{P}+\mathbf{P}\mathbf{A}=-\mathbf{Q}$，且 $\dot V=-\mathbf{x}^\top\mathbf{Q}\mathbf{x}$（取 $V=\mathbf{x}^\top\mathbf{P}\mathbf{x}$）。于是
\begin{equation}
  J=\int_{t_0}^{+\infty}\!\!-\dot V\,\mathrm dt=-V[\mathbf{x}(\infty)]+V[\mathbf{x}(t_0)]=\mathbf{x}^\top(t_0)\mathbf{P}\mathbf{x}(t_0),
  \label{eq:paramopt-closed}
\end{equation}
末步用渐近稳定的 $V[\mathbf{x}(\infty)]=0$。\textbf{这把一个无穷积分指标变成了闭式}：$J$ 只取决于初值 $\mathbf{x}(t_0)$ 与 $\mathbf{P}$，而 $\mathbf{P}$ 是 $\boldsymbol\xi$ 的函数。求 $J$ 对 $\boldsymbol\xi$ 的极小即得最优参数。

\begin{insight}[这正是 LQR 的“胚胎”]
\cref{eq:paramopt-closed} 中“代价 $=$ 初值的二次型 $\mathbf{x}_0^\top\mathbf{P}\mathbf{x}_0$、$\mathbf{P}$ 满足一个矩阵方程”这一结构，与导论 LQR（\cref{sec:ctrl-lqr}）的最优 cost-to-go $J^*(\mathbf{x})=\mathbf{x}^\top\mathbf{P}\mathbf{x}$ \emph{一模一样}——区别仅在于这里 $\mathbf{P}$ 满足\emph{李雅普诺夫}方程（给定反馈、只优化标量参数 $\boldsymbol\xi$），而 LQR 的 $\mathbf{P}$ 满足\emph{黎卡提}方程（同时优化整个反馈增益）。\textbf{从李氏方程到黎卡提方程，正是“参数最优化”到“最优控制”的跨越}——本章在此搭好了上游地基，严格的 CARE/DARE 求解与分离原理交 \cref{ch:lqr-lqg-riccati}。
\end{insight}

\begin{example}[阻尼比最优化]\label{ex:paramopt}
系统 $\dot{\mathbf{x}}=\big[\begin{smallmatrix}0&1\\-1&-2\xi\end{smallmatrix}\big]\mathbf{x}$，$\mathbf{x}(0)=(1,0)^\top$，代价 $\int_0^\infty\mathbf{x}^\top\mathbf{Q}\mathbf{x}\,\mathrm dt$，$\mathbf{Q}=\mathrm{diag}(1,\mu)$，$\mu>0$。求阻尼比 $\xi>0$ 使 $J$ 最小。由 $\mathbf{A}^\top\mathbf{P}+\mathbf{P}\mathbf{A}=-\mathbf{Q}$ 解得
\[
  \mathbf{P}=\begin{bmatrix}\xi+\frac{1+\mu}{4\xi}&\tfrac12\\[2pt]\tfrac12&\frac{1+\mu}{4\xi}\end{bmatrix},
  \qquad J=\mathbf{x}^\top(0)\mathbf{P}\mathbf{x}(0)=\xi+\frac{1+\mu}{4\xi}.
\]
令 $\mathrm dJ/\mathrm d\xi=0$ 得最优 $\xi=\tfrac{\sqrt{1+\mu}}{2}$。特例：$\mu=0$（指标退化为偏差平方积分 $\int_0^\infty y^2\mathrm dt$，$y=x_1$）时 $\xi=0.5$；$\mu=1$（指标 $\int_0^\infty(y^2+\dot y^2)\mathrm dt$）时 $\xi=0.707$——这正是经典控制中熟知的“最优阻尼比”结果，现在由李氏方程统一导出。
\end{example}

\subsection{用 \texorpdfstring{$V$}{V} 估计动态响应的快速性}
\label{sec:lyap-speed}

李氏函数 $V$ 是“到平衡的距离尺度”，$\dot V<0$ 表示状态趋向平衡的速度。定义\textbf{衰减率}
\begin{equation}
  \eta=-\frac{\dot V(\mathbf{x},t)}{V(\mathbf{x},t)},\qquad
  \eta_{\min}=\min_{\mathbf{x}}\Big[-\frac{\dot V}{V}\Big],
  \label{eq:eta-def}
\end{equation}
则 $\dot V\le-\eta_{\min}V$，积分得 $V(\mathbf{x},t)\le V(\mathbf{x}_0,t_0)\,e^{-\eta_{\min}(t-t_0)}$。$\eta_{\min}$ 越大，状态向原点收敛越快。若取 $V=\sum_i x_i^2$，则各分量满足 $|x_i(t)|\le|x_{i0}|\,e^{-\frac12\eta_{\min}(t-t_0)}$，对应经典意义下的\emph{时间常数}
\begin{equation}
  T=\frac{2}{\eta_{\min}}.
  \label{eq:time-const}
\end{equation}

对线性定常系统 $\dot{\mathbf{x}}=\mathbf{A}\mathbf{x}$，取 $V=\mathbf{x}^\top\mathbf{P}\mathbf{x}$、$\dot V=-\mathbf{x}^\top\mathbf{Q}\mathbf{x}$（$\mathbf{A}^\top\mathbf{P}+\mathbf{P}\mathbf{A}=-\mathbf{Q}$），则
\begin{equation}
  \eta_{\min}=\min_{\mathbf{x}}\frac{\mathbf{x}^\top\mathbf{Q}\mathbf{x}}{\mathbf{x}^\top\mathbf{P}\mathbf{x}}
  =\min\{\mathbf{x}^\top\mathbf{Q}\mathbf{x}:\mathbf{x}^\top\mathbf{P}\mathbf{x}=1\},
  \label{eq:eta-min-linear}
\end{equation}
这是一个广义瑞利商极小化，用拉格朗日乘子法（令 $\partial(\mathbf{x}^\top\mathbf{Q}\mathbf{x}+\mu(1-\mathbf{x}^\top\mathbf{P}\mathbf{x}))/\partial\mathbf{x}=2\mathbf{Q}\mathbf{x}-2\mu\mathbf{P}\mathbf{x}=0$）化为广义特征值问题 $\mathbf{Q}\mathbf{x}=\mu\mathbf{P}\mathbf{x}$，$\eta_{\min}$ 即最小广义特征值。

\begin{insight}[李氏函数给的是“保证下界”而非精确速率]
$\eta_{\min}$ 给出收敛快速性的一个\emph{保守估计}（下界）——它保证状态\emph{至少}以 $e^{-\frac12\eta_{\min}t}$ 衰减，但实际可能更快。这与 \cref{ch:dp-hjb} 中“HJB 值函数给最优代价的下界、任何可行李氏函数给一个上界估计”的对偶哲学相通：\emph{李雅普诺夫函数族里每一个都给一条保守界，取最紧的那条逼近真值}。这也呼应导论指数稳定定义（\cref{def:ctrl-exp-stable}）里的衰减率 $\beta$——$\eta_{\min}/2$ 就是一个可计算的 $\beta$ 下界。
\end{insight}

% ════════════════════════════════════════════════════════════════
\section*{本章常见误解汇总}
\label{sec:lyap-misconceptions}

\begin{table}[H]\centering\small
\caption{李雅普诺夫方法常见误解与正解}
\begin{tabular}{p{0.46\linewidth}p{0.46\linewidth}}
\toprule
\textbf{常见误解} & \textbf{正确理解} \\
\midrule
找不到李氏函数 $\Rightarrow$ 系统不稳定 & 判据是\emph{充分}条件；找不到只说明该工具失效，不能断言不稳定（\cref{thm:lyap-criteria}）\\
$\mathbf{A}$ 稳定 $\Rightarrow$ $\mathbf{A}+\mathbf{A}^\top\prec0$ & 错；须对给定 $\mathbf{Q}$ 解李氏方程得 $\mathbf{P}$（\cref{eq:lyap-P-int}），不能取 $\mathbf{P}=\mathbf{I}$ \\
矩阵对角元同号 $\Rightarrow$ 定号 & 须算\emph{顺序主子式}（\cref{thm:sylvester}）；对角元正未必正定 \\
线性化矩阵临界 $\Rightarrow$ 原系统临界稳定 & 临界（纯虚特征值）时第一法\emph{失效}，须用第二法/LaSalle（\cref{thm:lyap-indirect}）\\
传递函数极点全在左半平面 $\Rightarrow$ 系统稳定 & 那只是\emph{输出}稳定；零极点对消会掩盖内部不稳定（\cref{prop:output-stable}）\\
$\dot V$ 半负定 $\Rightarrow$ 只能判“稳定” & 若 $\dot V$ 沿轨迹不恒为零仍可判\emph{渐近}稳定（\cref{thm:lyap-criteria} + LaSalle）\\
李氏函数唯一 & 非唯一；不同 $V$ 影响判出的\emph{渐近稳定区域}大小（\cref{ex:vargrad}）\\
克拉索夫斯基法失败 $\Rightarrow$ 系统不稳定 & 其 $\mathbf{Q}(\mathbf{x})$ 处处正定条件过严；失败只表示该法无信息 \\
\bottomrule
\end{tabular}
\end{table}

% ════════════════════════════════════════════════════════════════
\section*{本章小结}
\label{sec:lyap-summary}

本章把导论“给会用”的李雅普诺夫理论\emph{展开为会推、会算、会造 $V$}：从 $\varepsilon$-$\delta$ 球域定义与能量直觉出发，建立第一法（线性化，临界失效）与第二法（能量法，四判据）；把线性李氏方程 $\mathbf{A}^\top\mathbf{P}+\mathbf{P}\mathbf{A}=-\mathbf{Q}$ 亲手推导并用于判稳定、定参数范围（连续/时变/离散三态）；攻克“怎样造 $V$”的难关（克拉索夫斯基法、变量梯度法，一题多解比较优劣）；最后补上 $V$ 的进阶用途（参数最优化、响应快速性估计）作为通向最优控制的桥。

\begin{table}[H]\centering\small
\caption{符号表}
\begin{tabular}{lll}
\toprule
符号 & 含义 & 出处 \\
\midrule
$\mathbf{x}_e$ & 平衡状态（约定平移到原点） & \cref{eq:equilibrium} \\
$S(\varepsilon),S(\delta)$ & 半径 $\varepsilon$/$\delta$ 的超球域 & \cref{def:lyap-isl} \\
$V(\mathbf{x})$ & 李雅普诺夫函数（广义能量/距离） & \cref{sec:lyap-second} \\
$\dot V(\mathbf{x})$ & $V$ 沿轨迹的时间导数（趋近速度） & \cref{thm:lyap-criteria} \\
$\mathbf{P},\mathbf{Q}$ & 二次型 $V=\mathbf{x}^\top\mathbf{P}\mathbf{x}$、$\dot V=-\mathbf{x}^\top\mathbf{Q}\mathbf{x}$ 的权重 & \cref{eq:lyap-eq} \\
$\Delta_i$ & $\mathbf{P}$ 的 $i$ 阶顺序主子式（Sylvester） & \cref{thm:sylvester} \\
$\mathbf{J}(\mathbf{x})$ & 雅可比 $\partial\mathbf{f}/\partial\mathbf{x}$（克拉索夫斯基法） & \cref{thm:krasovskii} \\
$\nabla V$ & $V$ 的梯度（变量梯度法待定量） & \cref{eq:vargrad-Vdot} \\
$\eta_{\min},T$ & 衰减率与时间常数 & \cref{eq:eta-def,eq:time-const} \\
\bottomrule
\end{tabular}
\end{table}

\begin{table}[H]\centering\small
\caption{定理与判据速查表}
\begin{tabular}{p{0.30\linewidth}p{0.40\linewidth}p{0.22\linewidth}}
\toprule
名称 & 一句话 & 出处 \\
\midrule
四种稳定性定义 & 稳定/渐近/大范围/不稳定的球域刻画 & \cref{def:lyap-isl,def:lyap-asymp} \\
第一法（间接法） & 看线性化矩阵特征值；临界失效 & \cref{thm:lyap-indirect} \\
第二法四判据 & $V$ 正定 + $\dot V$ 半负/负/正 $\to$ 稳定/渐近/不稳定 & \cref{thm:lyap-criteria} \\
希尔维斯特判据 & 顺序主子式判定号性 & \cref{thm:sylvester} \\
线性李氏方程 & 给 $\mathbf{Q}$ 解 $\mathbf{P}\succ0$ $\iff$ $\mathbf{A}$ Hurwitz & \cref{thm:lyap-eq-cont} \\
离散李氏方程 & $\mathbf{G}^\top\mathbf{P}\mathbf{G}-\mathbf{P}=-\mathbf{Q}$，谱半径 $<1$ & \cref{thm:lyap-disc} \\
时变离散李氏方程 & $\mathbf{G}^\top(k)\mathbf{P}(k{+}1)\mathbf{G}(k)-\mathbf{P}(k)=-\mathbf{Q}(k)$ & \cref{thm:lyap-tv-disc} \\
克拉索夫斯基法 & $V=\mathbf{f}^\top\mathbf{P}\mathbf{f}$，$\mathbf{J}^\top\mathbf{P}+\mathbf{P}\mathbf{J}\prec0$ & \cref{thm:krasovskii} \\
变量梯度法 & 猜 $\nabla V$、定系数、线积分还原 $V$ & \cref{algo:vargrad} \\
\bottomrule
\end{tabular}
\end{table}

% ════════════════════════════════════════════════════════════════
\section*{通向高级：本章交给后续什么}
\label{sec:lyap-bridge-up}

本章把李雅普诺夫方法做到了“会用、会推、会算、会造 $V$”的地基层；更严格、更前沿的内容由下列章节接手——它们\emph{默认}本章已立、只“回顾一行 + 交叉引用”：

\begin{itemize}
  \item \textbf{严格化与扩展} $\to$ \cref{ch:lyapunov-stability}（A–Q 的 D 章）：本章的四判据在那里获得完整的 $\varepsilon$-$\delta$ 严格证明；并向\emph{输入-状态稳定（ISS）}、\emph{反步法（backstepping）}、\emph{李雅普诺夫逆定理}（“稳定 $\Rightarrow$ 一定存在 $V$”——回答了本章“找不到 $V$ 怎么办”的存在性问题）、以及非自治系统的\emph{Barbalat 引理}（\cref{lem:ctrl-barbalat}）展开；
  \item \textbf{最优控制下界} $\to$ \cref{ch:dp-hjb}（B 章）：HJB 方程的值函数本质上是一个李雅普诺夫函数，“黏性解 + 值迭代/策略迭代”把本章 \cref{sec:lyap-speed} 的“保守界”哲学推向最优；
  \item \textbf{从李氏方程到黎卡提方程} $\to$ \cref{ch:lqr-lqg-riccati}（C 章）：本章 \cref{sec:lyap-paramopt} 的“代价 $=\mathbf{x}_0^\top\mathbf{P}\mathbf{x}_0$、$\mathbf{P}$ 满足矩阵方程”升级为 LQR 的 CARE/DARE，并给出分离原理的完整证明（本书不在 C0 重推 DARE/CARE，严格性交 C 章）；
  \item \textbf{无源性与 LaSalle} $\to$ 接 \cref{thm:ctrl-lasalle}：机器人 PD+重力补偿、力控的稳定性证明大量依赖“能量型 $V$ + LaSalle”，是本章能量直觉在多体系统上的落地；
  \item \textbf{安全对偶} $\to$ \cref{ch:clf-cbf}（F 章）：控制屏障函数（CBF）是李雅普诺夫函数的“安全版对偶”——后者保证\emph{收敛到}集合，前者保证\emph{不离开}集合。
\end{itemize}

\noindent 综合应用上，本章与 \cref{ch:eskf} 的卡尔曼滤波存在“估计-控制对偶”：李氏方程/黎卡提方程的形式同构于协方差传播（一句提示见 \cref{sec:lyap-tv-disc} 的 note），综合设计章（\cref{ch:linear-synthesis,ch:optimal-control-basics}）会把二者焊接成“观测器 + 状态反馈”的完整闭环。至此，“稳定性”这块地基已为整个控制部的收敛性论证备好统一语言。
